\documentclass[12pt,a4paper]{article}
\usepackage[utf8]{inputenc}
\usepackage[T1]{fontenc}
\usepackage[ngerman,latin]{babel}
\usepackage{amsmath,amssymb,amsthm}
\usepackage{geometry}
\usepackage{longtable}
\usepackage{booktabs}
\usepackage{array}
\usepackage{graphicx}
\usepackage{enumitem}
\usepackage{fancyhdr}
\usepackage{titlesec}
\usepackage{microtype}

\geometry{left=2.5cm,right=2.5cm,top=2.5cm,bottom=2.5cm}

\pagestyle{fancy}
\fancyhf{}
\fancyhead[C]{\small Carl Friedrich Gau\ss{} Werke --- Band 7}
\fancyfoot[C]{\thepage}
\renewcommand{\headrulewidth}{0.4pt}

\title{Carl Friedrich Gau\ss{} Werke --- Band 7, Teil 7}
\author{Carl Friedrich Gau\ss{}}
\date{}

\begin{document}

\maketitle

\section*{Nachtrag und Ver"offentlichungen}
\addcontentsline{toc}{section}{Nachtrag und Ver"offentlichungen}


\subsection*{[1. Zu Art. 7.]}
\textbf{Gen\"aherte Formeln f\"ur die gr\"osste Mittelpunktsgleichung.}

$e = \sin\varphi$

\medskip
\noindent
\textbf{I.} $[\varepsilon - M =]\ e + 2\arcsin\sqrt{(\sin\tfrac{1}{2}\varphi \tan\tfrac{1}{2}\varphi)}$, gibt: $2e + \tfrac{1}{4}e^3 + \tfrac{1}{96}e^5$, Fehler: $-\tfrac{107}{2880}e^7$

\medskip
\noindent
\textbf{II.} $[\varepsilon - M =]\ e + 2\arcsin\sqrt{(\sin\tfrac{1}{2}\varphi \tan\tfrac{1}{2}\varphi)}$, gibt: $2e + \tfrac{1}{4}e^3 + \tfrac{15}{96}e^5$, Fehler: $-\tfrac{107}{2880}e^7$

\newpage

\subsection*{[2. Zu Art. 53.]}
\textbf{Projection einer Planetenbahn auf die Ekliptik.}

$e$, $\varphi$, $a$, $i$, $\Omega$, $E$ in gew\"ohnlicher Bedeutung

$\omega$ Perihel

$\lambda$ hel[iocentrische] L[\"a]nge

$p$ curtirter Abstand [von der Sonne].

[F\"ur die rechtwinkligen Coordinaten in der Ebene der Ekliptik, gez\"ahlt
von der Knotenlinie als Abscissenlinie hat man nach Art.~53 und 58 der
\textit{Theoria motus corporum coelestium}:]
\begin{align*}
a\cos\psi(\cos E - e) - a\cos\varphi\sin\psi\sin E &= p\cos(\lambda - \Omega) \\
a\cos\psi\sin\psi(\cos E - e) + a\cos\varphi\cos\psi\cos i\sin E &= p\sin(\lambda - \Omega).
\end{align*}

[Man setze
\begin{align*}
a\cos\psi &= \alpha\cos\varepsilon\cos\gamma - \beta\sin\varepsilon\sin\gamma \\
a\cos\varphi\sin\psi &= \alpha\sin\varepsilon\cos\gamma + \beta\cos\varepsilon\sin\gamma \\
a\cos\psi\sin i &= \alpha\cos\varepsilon\sin\gamma + \beta\sin\varepsilon\cos\gamma \\
a\cos\varphi\cos i\cos\psi &= -\alpha\sin\varepsilon\sin\gamma + \beta\cos\varepsilon\cos\gamma,]
\end{align*}

\newpage

\noindent
man bestimme [also] $\alpha$, $\beta$, $\gamma$, $\varepsilon$ durch die Gleichungen
\begin{align*}
a(1 - \cos\psi\cos i)\cos\gamma &= (\alpha - \beta)\cos(\gamma - \varepsilon) \\
a(\cos i - \cos\psi)\sin\gamma &= (\alpha - \beta)\sin(\gamma - \varepsilon) \\
a(1 + \cos\psi\cos i)\cos\gamma &= (\alpha + \beta)\cos(\gamma + \varepsilon) \\
a(\cos i + \cos\psi)\sin\gamma &= (\alpha + \beta)\sin(\gamma + \varepsilon),
\end{align*}
[so wird
\begin{align*}
a(\cos(E + \varepsilon) - e\cos\varepsilon)\cos\gamma - \beta(\sin(E + \varepsilon) - e\sin\varepsilon)\sin\gamma &= p\cos(\lambda - \Omega) \\
a(\cos(E + \varepsilon) - e\cos\varepsilon)\sin\gamma + \beta(\sin(E + \varepsilon) - e\sin\varepsilon)\cos\gamma &= p\sin(\lambda - \Omega)
\end{align*}
oder]
\begin{align*}
\alpha\cos(E + \varepsilon) - \alpha e\cos\varepsilon &= p\cos(\lambda - \gamma - \Omega) \\
\beta\sin(E + \varepsilon) - \beta e\sin\varepsilon &= p\sin(\lambda - \gamma - \Omega).
\end{align*}

Hier sind $\alpha$, $\beta$ die beiden halben Hauptaxen der Projectionsellipse;
$-\alpha e\cos\varepsilon$, $-\beta e\sin\varepsilon$ die Coordinaten des Mittelpunkts, [$p\cos(\lambda - \gamma - \Omega)$,
$p\sin(\lambda - \gamma - \Omega)$ die rechtwinkligen Coordinaten des Planeten,] die [durch die
Sonne zur] Axe $2\alpha$ [gezogene Parallele] als Abscissenlinie angenommen; $\gamma + \Omega$
deren Neigung gegen die \"Aquinoktiallinie.

\medskip
\noindent
Bei der Berechnung kann man noch folgende H\"ulfsgr\"ossen einf\"uhren:
\[
\cos i = \tan k, \quad \cos\varphi = \tan l, \quad \frac{a}{\cos k \cos i} = c;
\]
alsdann ist
\begin{align*}
a(1 - \cos\psi\cos i) &= c\cos(k + l) \\
a(\cos i - \cos\psi) &= c\sin(k - l) \\
a(1 + \cos\psi\cos i) &= c\cos(k - l) \\
a(\cos i + \cos\psi) &= c\sin(k + l).
\end{align*}

\newpage

\textbf{Beispiel. Juno. Elemente f\"ur 1810.}

\begin{center}
\begin{tabular}{ll}
$\psi = 14^\circ 45' 35''$ & \\
$\log a = 0{,}42648$ & mot.~diurn.~trop. $= 813{,}5165$ \quad sid. $= 813{,}379$ \\
$i = 13^\circ 4' 12''$ & \\
$\beta = 171^\circ 7' 48''$ & \\
$\omega = 53^\circ 10' 52''$ & Epoche [der mittl.~L\"ange] 1810 $= 95^\circ 33' 50''$, \quad $\varepsilon = 242^\circ 3' 4''$
\end{tabular}
\end{center}

\begin{center}
\begin{tabular}{ll}
$\cos i \ldots\ldots 9{,}988\ 6011$ & \\
$\cos\varphi \ldots\ldots 9{,}985\ 4277$ & \\
$k = 44^\circ 14' 53{,}40''$ & $l = 44^\circ 2' 20{,}34''$ \\
$\cos\psi \ldots\ldots 9{,}855\ 1097$ & \\
$\cos l \ldots\ldots 9{,}856\ 6489$ & \\
\multicolumn{2}{l}{$9{,}711\ 7586$} \\
$0{,}42648$ & \\
$c \ldots\ldots 0{,}714\ 72$ & \\
$k + l = 88^\circ 17' 13{,}74''$ & \\
$k - l = 0^\circ 12' 33{,}26''$ & \\
$\cos(k + l) \ldots 8{,}47555$ & $\sin(k + l) \ldots 9{,}99981$ \\
$\cos(k - l) \ldots 0{,}00000$ & $\sin(k - l) \ldots 7{,}56252$ \\
\multicolumn{2}{l}{$9{,}67088$} \\
$\sin\psi \ldots\ldots 9{,}946\ 14$ & \\
\end{tabular}
\end{center}

\begin{center}
\begin{tabular}{ll}
$(\alpha - \beta)\cos(\gamma - \varepsilon) \ldots 8{,}861\ 15$ & $(\alpha - \beta)\sin(\gamma - \varepsilon) \ldots 8{,}223\ 38$ \\
$(\alpha + \beta)\cos(\gamma + \varepsilon) \ldots 0{,}385\ 60$ & $(\alpha + \beta)\sin(\gamma + \varepsilon) \ldots 0{,}660\ 67$
\end{tabular}
\end{center}

\begin{center}
\begin{tabular}{l}
$= \begin{cases} 192^\circ 58' 2'' \\ 242^\circ 2' 26'' \end{cases}$ \\
$\alpha - \beta \ldots\ldots 8{,}872\ 37$ \\
$\alpha + \beta \ldots\ldots 0{,}714\ 57$
\end{tabular}
\end{center}

\begin{center}
\begin{tabular}{ll}
$\gamma = 21^\circ 7' 30{,}14''$ & \\
$\varepsilon = 24^\circ 32' 12''$ & \\
$\gamma + \Omega = 28^\circ 38' 2''$ & \\
$\alpha \ldots\ldots 0{,}419\ 74$ & \\
$e \ldots\ldots 0{,}160\ 11$ & \\
$\alpha e \ldots\ldots 9{,}825\ 88$ & \\
$\beta e \ldots\ldots 9{,}958\ 90$ & \\
$\beta e \ldots\ldots 9{,}813\ 39$ & \\
$\sin\varepsilon \ldots 9{,}618\ 34$ & \\
$-\alpha e\cos\varepsilon = -0{,}609\ 23$ & Coord. \\
$-\beta e\sin\varepsilon = -0{,}270\ 22$ & Centri.
\end{tabular}
\end{center}

[Man kann also die heliocentrischen Coordinaten jetzt nach der Formel
\begin{align*}
p\cos(\lambda - 28^\circ 38' 2'') &= -0{,}60923 + 0{,}41974 \cdot \cos(E + 24^\circ 32' 12'') \\
p\sin(\lambda - 28^\circ 38' 2'') &= -0{,}27022 + 0{,}40725 \cdot \sin(E + 24^\circ 32' 12'')
\end{align*}
berechnen, sobald nur die excentrische Anomalie bekannt ist.

Die Zahlen in Parenthese bedeuten Logarithmen.]

Die Hauptmomente der Aufl\"osung der umgekehrten Aufgabe: aus der
Projection einer Planetenbahn und ihres Brennpunkts, jene selbst zu finden,
sind folgende:

\newpage

Es sei $\alpha$ die halbe grosse, $\beta$ die halbe kleine Axe der Projectionsellipse;
ferner, vom Mittelpunkt der Ellipse an gerechnet, gegen eine willk\"urliche
Abscissenlinie

$P \ldots$ Richtung von $\alpha$

$Q \ldots$ Richtung der Linie zum Brennpunkte, deren Gr\"osse $= \delta$.

Man setze
\begin{equation}\label{eq:6}
\tag{6}
\tan(R - P) \cdot \tan(Q - P) = -\frac{\beta\beta}{\alpha\alpha},
\end{equation}
[so dass also $R$ die Richtung des zur Linie zum Brennpunkt conjugirten
Durchmessers ist; wenn $q$ den durch den Brennpunkt gehenden Halbmesser
und $r$ den ihm conjugirten bezeichnet, so ist
\[
\frac{r}{q} = \frac{\sqrt{(\alpha\alpha + \delta\delta)}}{\alpha},
\]
also
\begin{equation}\label{eq:7}
\tag{7}
\frac{r\delta}{q\alpha} = \frac{\sin 2(Q - P)}{\sin 2(R - P)},
\end{equation}
[Die Excentricit\"at der Ellipse ist]
\begin{equation}\label{eq:9}
\tag{9}
\frac{\delta}{\alpha} = e.
\end{equation}

[Durch Einf\"uhrung des auch im vorigen benutzten Winkels $\varepsilon$, welcher
die zum Radius Vector $q$ geh\"orige excentrische Anomalie in der Projections-
ellipse darstellt, lassen sich die Endpunkte der conjugirten Radii Vectores $q$
und $r$ in rechtwinkligen Coordinaten, jetzt bezogen auf die Axe $2\alpha$ als Ab-
scissenlinie, wie folgt, ausdr\"ucken:
\begin{align*}
q\cos(Q - P) &= \alpha\cos\varepsilon, & r\cos(R - P) &= -\alpha\sin\varepsilon, \\
q\sin(Q - P) &= \beta\sin\varepsilon, & r\sin(R - P) &= \beta\cos\varepsilon;
\end{align*}
leitet man aus diesen die rechtwinkligen Coordinaten, bezogen auf die Knoten-
linie als Abscissenlinie ab, indem man um den Winkel $\gamma = P - \Omega$ dreht,
so hat man mit R\"ucksicht auf die Gleichungen 2:]
\begin{align*}
x\cos(Q - \Omega) &= \alpha\cos\varepsilon, & -x\sin(R - \Omega) &= -\alpha\sin\varepsilon \\
y\sin(Q - \Omega) &= \alpha\cos i\sin\psi, & y\cos(R - \Omega) &= \alpha\cos\psi\cos\varepsilon.
\end{align*}

\newpage

[woraus man leicht ableitet]
\begin{equation}\label{eq:11}
\tag{11}
\frac{rr}{qq} = \frac{\alpha\alpha(1 - \cos i)}{\beta\beta(1 + \cos i)} \cdot \frac{\sin^2(Q - \Omega)}{\sin^2(R - \Omega)}
\end{equation}

[Ferner findet sich]
\begin{equation}\label{eq:12}
\tag{12}
\frac{rr}{qq} = \frac{\alpha\alpha(1 - \cos i)}{\beta\beta(1 + \cos i)}
\end{equation}

\begin{align*}
qq + rr &= \alpha\alpha(1 + \cos i) \\
\sin(R - Q) &= \alpha\alpha\cos i \\
\sqrt{(1 - ee)} \cdot \cos(R - Q) &= -\alpha\alpha\sin i \cdot \sin 2\psi \\
qq - rr &= \alpha\alpha\sin i \cdot \cos 2\psi.
\end{align*}

[In \"ahnlicher Weise leitet man ab
\begin{align*}
-rr\cos(Q + R - 2\Omega) &= -\tfrac{1}{2}\alpha\alpha(1 + \cos i)\sin 2\psi \\
rr\sin(Q + R - 2\Omega) &= \alpha\alpha\cos i \cdot \cos 2\psi,
\end{align*}
also mit R\"ucksicht auf die Gleichungen 12:]
\begin{equation}\label{eq:13}
\tag{13}
\begin{aligned}
\alpha\alpha\sin i \cdot \cos(Q + R - 2\Omega) &= (qq + rr)\cos(R - Q) \\
\alpha\alpha\sin i \cdot \sin(Q + R - 2\Omega) &= (qq - rr)\sin(R - Q).
\end{aligned}
\end{equation}

[Die abgeleiteten Gleichungen enthalten alles n\"othige zur L\"osung der
Aufgabe und k\"onnen auf mannigfaltige Weise benutzt werden; sind $\alpha$, $\beta$, $P$,
$Q$, $\delta$ gegeben, so kann man z.~B. $R$ aus 6, $q$ und $r$ aus 7, $e$ aus 9 bestimmen;
sodann finden sich $\Omega$, $a$, $i$, $\psi$ aus 13 und 10. Die Lage der Abscissenlinie,
von der die Richtungen $P$, $Q$, $R$, $\psi$ gerechnet werden, bleibt vollkommen will-
k\"urlich.]

\subsection*{[3. Zu Art. 70.]}

Da in den Gleichungen Art.~62 der Winkel $N$ willk\"urlich ist, so setze
man $N = \lambda$; dadurch werden jene, wegen $r' = r\cos\beta$, $R' = R\cos B$, $A' = A\cos b$:
\begin{align}
\tag{1} r\cos\beta - R\cos B \cos(L - \lambda) &= A\cos b \cos(l - \lambda) \\
\tag{2} R\cos B \sin(L - \lambda) &= A\cos b \sin(l - \lambda) \\
\tag{3} r\sin\beta - R\sin B &= A\sin b.
\end{align}

\newpage

Aus 2 und 3 folgt
\[
\tan(l - \lambda) = \frac{R\cos B \sin(L - \lambda)}{r\cos\beta - R\cos B \cos(L - \lambda)}
\]
oder n\"aherungsweise
\begin{equation}\tag{4}
\tan(l - \lambda) = \frac{R\cos B \sin(L - \lambda)}{r\cos\beta}.
\end{equation}

Ferner folgt aus 1 und 3
\[
R\cos B \sin\beta \cos(L - \lambda) - R\sin B \cos\beta = A(\sin b \cos\beta - \cos b \sin\beta \cos(l - \lambda))
\]
oder
\begin{equation}\tag{5}
\frac{R\cos B \cos\beta(\tan\beta - \cos(L - \lambda)\tan B)}{\sin\beta} = \frac{A\sin(b - \beta)}{\sin(l - \lambda)}
\end{equation}
oder n\"aherungsweise

\subsection*{[4. Zu Art. 76--77.]}
\textbf{Allgemeine Differentialformeln f\"ur die geocentrischen \"Orter der Planeten.}

$l$ geocentrische L\"ange des Planeten, $l - \Omega = \nu$

$L$ heliocentrische L\"ange der Erde

$u$ Argument der Breite

$b$ geocentrische Breite

$v$ wahre Anomalie

$\omega$ L\"ange des Perihels

$n$ mittlere t\"agliche Bewegung

$\text{Ep.}$ mittlere L\"ange f\"ur die Epoche

$\tau$ seit der Epoche verflossene Zeit in Tagen.

\begin{align*}
\text{1. } & a\tan i \sin u = A = \left(\frac{dr}{d\Omega}\right) &
\text{3. } & \frac{aa\cos v}{r} \cdot \beta = B = \left(\frac{dl}{d\omega}\right) \\
\text{2. } & -a\cos u \cos v = \alpha = \left(\frac{dr}{di}\right) &
\text{4. } & -\frac{aa\sin E}{r} \cdot \frac{2 + e\cos E}{\sqrt{(1 - ee)}} = \beta = \left(\frac{dl}{de}\right)
\end{align*}

\begin{align*}
\text{5. } & \tan\zeta = \frac{r}{R\cos(L - l)} \\
\text{6. } & \sin\zeta \tan b = \tan N; \quad \cos N = \frac{\cos i}{\cos b}, \quad \sin N = \sin M \sin i \\
\text{7. } & \frac{\sin\zeta \sin(L - l)}{\cos\zeta \sin u} = \tan P \\
\text{8. } & -\sin N \cdot \beta = x
\end{align*}

\newpage

\begin{align*}
\text{9. } & \frac{\sin\zeta \cos b}{r} = m \\
\text{10. } & -\lambda \cot(M - u) = B = -\frac{\beta}{m} \\
\text{11. } & -\cos N \tan b = C = \frac{\gamma}{m} \\
\text{12. } & \frac{\cos(\zeta - \nu)}{R} = y \\
\text{13. } & x + y = \frac{1}{r} = (w) \\
\text{14. } & v = \frac{1}{r} = m \\
\text{15. } & \frac{r\sin M \cos i \cos(N - b) - P}{A\sin N} = \frac{\delta}{m} \\
\text{16. } & \frac{r\sin M \cos i \cos(N - b)}{A\cos N} = \frac{\varepsilon}{m} \\
\text{17. } & \alpha\sin b \cos b = D' \\
\text{18. } & A\alpha + B\beta = E \\
\text{19. } & A'\alpha + B'\beta = E' \\
\text{20. } & A\alpha' + B\beta' = F \\
\text{21. } & A'\alpha' + B'\beta' = F'
\end{align*}

\begin{align*}
dl &= E\,d\text{Ep.} + (nE - \tfrac{\delta}{m})\,d\tau + (B - E)\,d\Omega + F\,d\varphi + C\,di + (D - E)\,d\Omega \\
db &= E'\,d\text{Ep.} + (nE' - \tfrac{\varepsilon}{m})\,d\tau + (B' - E')\,d\Omega + F'\,d\varphi + C'\,di + (D' - E')\,d\Omega.
\end{align*}

\subsection*{[5. Zu Art. 78 II.]}
\textbf{Direkte Aufl\"osung der Gleichungen}
\begin{align}
\tag{1} \sin(P - A) &= k\sin(Q - B) \\
\tag{2} \sin(P - A') &= k'\sin(Q - B)
\end{align}
[wo $A$, $A'$, $B$, $B'$, $k$, $k'$ gegebene Gr\"ossen und $P$, $Q$ zu bestimmen sind].

Es folgt daraus
\begin{align*}
\tan\tfrac{1}{2}(P + Q - A - B) &= \frac{k - 1}{k + 1} \tan\tfrac{1}{2}(P - Q - A + B) \\
\tan\tfrac{1}{2}(P + Q - A' - B') &= \frac{k' - 1}{k' + 1} \tan\tfrac{1}{2}(P - Q - A' + B')
\end{align*}
oder wenn wir
\begin{align*}
\tfrac{1}{2}(P + Q - A - B) &= x, & \tfrac{1}{2}(P - Q - A + B) &= y \\
\tfrac{1}{2}(A + B - A' - B') &= \alpha, & \tfrac{1}{2}(A - B - A' + B') &= \beta \\
m &= \frac{1 - k}{1 + k}, & n &= \frac{1 - k'}{1 + k'}
\end{align*}
setzen:
\begin{align}
\tag{3} \tan x &= m\tan y \\
\tag{4} \tan(x + \alpha) &= n\tan(y + \beta)
\end{align}

\newpage

oder in imagin\"arer Form
\[
\frac{e^{ix} - e^{-ix}}{e^{ix} + e^{-ix}} = m \cdot \frac{e^{iy} - e^{-iy}}{e^{iy} + e^{-iy}}
\]
[und analog f\"ur Gleichung 4] oder
\[
\frac{e^{2ix} - 1}{e^{2ix} + 1} = m \cdot \frac{e^{2iy} - 1}{e^{2iy} + 1}
\]
und
\[
\frac{e^{2i(x+\alpha)} - 1}{e^{2i(x+\alpha)} + 1} = n \cdot \frac{e^{2i(y+\beta)} - 1}{e^{2i(y+\beta)} + 1}
\]
folglich
\begin{multline*}
((1 + m)e^{iy} + (1 - m)e^{-iy})((1 - n)e^{i(x+\alpha)} + (1 + n)e^{-i(x+\alpha)})e^{i\beta} \\
= ((1 - m)e^{iy} + (1 + m)e^{-iy})((1 + n)e^{i(x+\alpha)} + (1 - n)e^{-i(x+\alpha)})e^{-i\beta}
\end{multline*}
oder wenn man entwickelt und mit 2 dividirt:
\begin{multline*}
e^{i(x-y+\alpha+\beta)}(m - n\cos\alpha + i(1 - mn)\sin\alpha) \\
- e^{-i(x-y+\alpha+\beta)}(m - n\cos\alpha - i(1 - mn)\sin\alpha) \\
= 2i(m + n\cos\alpha \sin\beta - (1 + mn)\sin\alpha \cos\beta)
\end{multline*}

F\"uhrt man also die H\"ulfsgr\"ossen $g$, $G$, $h$, $H$ ein, so dass
\begin{align*}
(m - n)\cos\alpha &= g\cos G, & (m + n)\cos\alpha &= h\cos H \\
(1 - mn)\sin\alpha &= g\sin G, & (1 + mn)\sin\alpha &= h\sin H,
\end{align*}
so wird diese Gleichung
\begin{equation}\tag{5}
g\sin(2y + \beta - G) = h\sin(\beta - H),
\end{equation}
woraus die beiden Werthe von $y$ sich ergeben. Die correspondirenden Werthe
von $x$ finden sich dann aus (3) oder (4).

Man kann auch setzen
\begin{align*}
(k' - k)\cos\alpha &= g'\cos G, & (kk' - 1)\cos\alpha &= h'\cos H \\
-(k' + k)\sin\alpha &= g'\sin G, & (kk' + 1)\sin\alpha &= h'\sin H
\end{align*}
und hat dann
\[
g'\sin(2y + \beta - G) = h'\sin(\beta - H).
\]
[\"Ubrigens ist hier
\[
g^2 = (k - k')(k' - k), \quad g'^2 = (k - k')(k' - k).
\]
$h^2 = (k - k')(k' - k)h$.]

\subsection*{[6. Zu Art. 88--90.]}

Um aus $\sin\tfrac{1}{2}\varphi$ den Logarithmen von $\sin\varphi$ zu finden, bediene man sich
folgender N\"aherungsformel:
\[
\log\sin\varphi = \log(2\sin\tfrac{1}{2}\varphi) + \tfrac{1}{2}\mu(\sin\tfrac{1}{2}\varphi)^2 + \tfrac{1}{2}\mu(\sin\tfrac{1}{2}\varphi)^4 + \tfrac{1}{3}\mu(\sin\tfrac{1}{2}\varphi)^6
\]
Hiedurch findet man den gesuchten Logarithmen zu gross: folgende Tafel
gibt die abzuziehende Correction in der $7^{\text{ten}}$ Decimale an.

\begin{center}
\begin{tabular}{|c|c|}
\hline
Grenze & Abzuziehende \\
$\log\sin\tfrac{1}{2}\varphi$ & Correction \\
\hline
9,069 & 0 \\
9,147 & 26'' 54' \\
9,184 & 32 18 \\
9,208 & 35 9 \\
9,226 & 37 10 \\
9,240 & 38 45 \\
9,252 & 40 3 \\
9,262 & 41 10 \\
9,271 & 42 10 \\
9,279 & 43 3 \\
9,286 & 43 51 \\
9,293 & 44 35 \\
\hline
\end{tabular}
\end{center}

\newpage

\textbf{Vorschriften, um den Logar.~Sinus eines kleinen Bogens zu finden.}

$\log\sin n'' = \log\sin\varphi = \log\varphi - M(\tfrac{1}{6}\varphi\varphi + \tfrac{1}{180}\varphi^4 + \tfrac{1}{2835}\varphi^6 + \tfrac{1}{37800}\varphi^8 + \text{etc.})$

Man bilde die Gr\"ossen $\lambda$, $\lambda'$, $\lambda''$, $\mu$, $\mu'$, $\mu''$, $\ldots$
\begin{align*}
\lambda &= 2\log n + 8{,}230\ 7828 \\
\lambda' &= \lambda + \tfrac{1}{6}\mu(\tfrac{1}{2}\lambda' - \tfrac{1}{2}\lambda) \\
\lambda'' &= \lambda' + \tfrac{1}{6}\mu'(\tfrac{1}{2}\lambda'' - \tfrac{1}{2}\lambda') \\
\lambda''' &= \lambda'' + \tfrac{1}{6}\mu''(\tfrac{1}{2}\lambda''' - \tfrac{1}{2}\lambda'') \\
&\text{u.~s.~w.}
\end{align*}
nach folgendem Gesetze:
\begin{center}
\begin{tabular}{ll}
$\log\mu$ & $\log\mu'$ \\
$\log\mu'$ & $\log\mu''$ \\
$\log\mu''$ & $\log\mu'''$ \\
$\vdots$ & $\vdots$
\end{tabular}
\end{center}

\begin{center}
$38^*$
\end{center}

so ist
\[
\log\sin\varphi = 4{,}685\ 5749 - 10 + \log n - \tfrac{1}{6}\mu'''
\]

\textbf{Beispiel:} $\varphi = 37^\circ 6' = 133560''$.

\begin{center}
\begin{tabular}{r|l}
$4{,}685\ 5749$ & $8{,}230\ 7828$ \\
$\log n \quad 5{,}125\ 6764$ & $2\log n \ldots 10{,}251\ 3528$ \\
\hline
$9{,}811\ 2513$ & $\lambda = 8{,}482\ 1356$ \quad $\mu = 0{,}030\ 3483{,}85$ \\
& $+ 60696{,}8$ \quad $4271{,}25$ \\
$\mu''' = 0{,}030\ 7842$ & $\lambda' = 8{,}488\ 2052{,}8$ \quad $\mu' = 0{,}030\ 7755{,}1$ \\
$+ 1200{,}0$ \quad $85{,}0$ & \\
$\lambda'' = 8{,}488\ 3252{,}8$ & $\mu'' = 0{,}030\ 7840{,}1$ \\
$+ 24{,}3$ & \\
$\lambda''' = 8{,}488\ 3277{,}1$ & $\mu''' = 0{,}030\ 7841{,}9$.
\end{tabular}
\end{center}

\begin{center}
$\log\sin\varphi \ldots 9{,}780\ 4671$
\end{center}

\subsection*{[7. Zu Art. 90 und 100.]}

\textit{Astronomisches Jahrbuch f\"ur 1814, S.~256. Berlin 1811.}

Noch f\"uge ich Ihrem Wunsche zufolge einen kleinen Zusatz zu
meiner \textit{Theoria motus corporum coelestium} bei.

Zur Aufl\"osung der wichtigen Aufgabe, aus zweien Radiis vectoribus und
dem eingeschlossenen Winkel die elliptischen oder hyperbolischen Elemente
zu bestimmen, habe ich mich mit grossem Vortheil einer H\"ulfsgr\"osse $\Phi$ bei
der Ellipse, $C$ bei der Hyperbel bedient, f\"ur welche ich jenem Werke eine
Tafel angeh\"angt habe. Berechnet ist diese Tafel nach einem dort angef\"uhrten
continuirlichen Bruche, dessen vollst\"andige Ableitung aber dort nicht gegeben
ist, und zu dessen theoretischer Entwickelung, die mit andern Untersuchungen
zusammenh\"angt, ich bisher noch nicht Gelegenheit gefunden habe. Es wird
daher manchem lieb sein, hier einen andern Weg angezeigt zu finden, auf
welchem man jene H\"ulfsgr\"osse eben so bequem h\"atte berechnen k\"onnen.

Wir haben (Art.~90):
\[
\Phi = \frac{X - \tfrac{1}{2}X^2 + \tfrac{1}{3}X^3 - \ldots}{1 - \tfrac{1}{2}X + \tfrac{1}{3}X^2 - \ldots}
\]

\newpage

Der Z\"ahler dieses Bruchs verwandelt sich leicht, wenn man f\"ur $X$ die
dort gegebene Reihe substituirt, in
\[
\tfrac{1}{2}x\left(1 + \tfrac{2 \cdot 8}{3 \cdot 9}x + \tfrac{3 \cdot 8 \cdot 10}{3 \cdot 9 \cdot 11}x^2 + \tfrac{4 \cdot 8 \cdot 10 \cdot 12}{3 \cdot 9 \cdot 11 \cdot 13}x^3 + \tfrac{5 \cdot 8 \cdot 10 \cdot 12 \cdot 14}{3 \cdot 9 \cdot 11 \cdot 13 \cdot 15}x^4 + \text{etc.}\right).
\]
Setzt man also die Reihe
\[
A = 1 + \tfrac{2 \cdot 8}{3 \cdot 9}x + \tfrac{3 \cdot 8 \cdot 10}{3 \cdot 9 \cdot 11}x^2 + \tfrac{4 \cdot 8 \cdot 10 \cdot 12}{3 \cdot 9 \cdot 11 \cdot 13}x^3 + \text{etc.},
\]
so wird
\[
\Phi X - \tfrac{1}{2}X^2 + \tfrac{1}{3}X^3 - \ldots = \tfrac{1}{2}x \cdot A,
\]
\[
\Phi = \frac{\tfrac{1}{2}x \cdot A}{1 - \tfrac{1}{2}X + \tfrac{1}{3}X^2 - \ldots} = \frac{\tfrac{1}{2}x \cdot A}{1 - \tfrac{1}{2}A x},
\]
nach welcher Formel man $\Phi$ immer bequem und sicher berechnen kann. F\"ur
$C$ braucht man nur $z$ anstatt $x$ zu setzen.

Ich bemerke nur noch, dass man $A$ noch bequemer nach folgender Formel
berechnen kann
\[
A = \left(\tfrac{1}{1 - \tfrac{8}{9}x}\right)^2 \cdot \left(1 + \tfrac{1}{3}\cdot\tfrac{8}{9}x + \tfrac{1 \cdot 4}{3 \cdot 6}\left(\tfrac{8}{9}x\right)^2 + \tfrac{1 \cdot 4 \cdot 7}{3 \cdot 6 \cdot 9}\left(\tfrac{8}{9}x\right)^3 + \text{etc.}\right),
\]
allein die Ableitung dieser Reihe aus der vorigen beruht auf Gr\"unden, die hier
nicht ausgef\"uhrt werden k\"onnen [$^*$].

\begin{flushright}
[$^*$ Vgl. Band III, S.~208.]
\end{flushright}

\subsection*{[8. Zu Art. 114 und 177.]}

\textit{Druckfehler in Dr.~Gauss' Theoria motus corporum coelestium etc.\ Hamburgi 1809, vom Senator Bar.~Oriani.}

\textit{Monatliche Correspondenz Band XXI, S.~282--283, 1810 M\"arz.}

\ldots Ut aequatio 7 [art.~114] fiat identica, debent coefficientes ipsorum
$\delta\delta'\delta''$, $\delta\delta'\delta''$, $\delta\delta'\delta''$ etc.~esse $= 0$; sed post debitas reductiones coefficiens
ipsius $\delta\delta'\delta''$ est
\begin{multline*}
\tan\beta' \tan\beta'' \sin(L'' - a'') \sin(L' - a) + \tan\beta'' \tan\beta \sin(L' - a') \sin(L - a'') \\
+ \tan\beta \tan\beta' \sin(L - a) \sin(L'' - a').
\end{multline*}

Hinc si habeatur tantummodo
\begin{multline*}
\tan\beta' \tan\beta'' \sin(L'' - a'') \sin(L' - a) \\
+ \tan\beta'' \tan\beta \sin(L' - a') \sin(L - a'') \\
+ \tan\beta \tan\beta' \sin(L - a) \sin(L'' - a') = 0,
\end{multline*}
idem coefficiens non fit $= 0$, sed requiritur praeterea ut sit
\[
B = 0, \quad B' = 0, \quad B'' = 0.
\]

\textit{Elegans theorema, quod tribuitur [art.~117] illustr.~Laplace, revera a
Leonardo Eulero primum inventum est. Et enim in Comment.~Acad.~Petrop.
Tom.~XVI [p.~118] Eulero ostendit, integrale $\int \frac{dx}{\sqrt{(-\log x)}}$ sumtum ab $x = 1$
ad $x = 0$ esse $= \sqrt{\pi}$, existente $\pi$ semicircumferentia circuli, radio $= 1$ de-
scripti. Iamvero ponendo $x = e^{-tt}$ habetur
\[
\int \frac{dx}{\sqrt{(-\log x)}} = \int \frac{-2te^{-tt}\,dt}{t} = 2\int e^{-tt}\,dt.
\]
Ideoque integrale $\int e^{-tt}\,dt$ a $t = 0$ ad $t = \infty$ erit $= \tfrac{1}{2}\sqrt{\pi}$, et propterea
idem integrale a $t = -\infty$ ad $t = +\infty$ fiet $= \sqrt{\pi}$.]}

\begin{flushright}
[$^*$ Handschriftliche Bemerkung von Gau\ss:] Dies Theorem findet sich a.a.O.~nicht, wohl aber
p.~101 folgendes: $\int \sqrt{\log\frac{1}{x}}\,dx = \tfrac{1}{2}\sqrt{\pi}$. Schreibt man hier $x = e^{-tt}$, so wird
$\int 2tte^{-tt}\,dt = \tfrac{1}{2}\sqrt{\pi}$. Es ist aber $\int 2tte^{-tt}\,dt = -\int d(te^{-tt}) + \int e^{-tt}\,dt$ und
$te^{-tt} = 0$ sowohl f\"ur $t = 0$ als f\"ur $t = \infty$, also $\int_0^\infty e^{-tt}\,dt = \tfrac{1}{2}\sqrt{\pi}$.
\end{flushright}

\newpage

\subsection*{[9. Zu Art. 138.]}

Das Criterium des Falls, wo der die geocentrische Bewegung tangirende
gr\"osste Kreis durch den Sonnenort geht, l\"asst sich sehr elegant auf die helio-
centrischen Positionen beziehen.

Es seien

$R$, $r$ Abst\"ande von der Sonne f\"ur Erde und Planet;

$L$, $l$ L\"angen der Erde und des Planeten, letztere in der Bahn und
beide vom aufsteigenden Knoten gez\"ahlt;

$E$, $e$ Excentricit\"aten;

$V$, $v$ wahre Anomalien;

$P$, $p$ halbe Parameter;

dann ist das Criterium jenes Falls in der Gleichung
\[
\frac{\sqrt{p} \cdot \sin l}{1 + e\cos v} = \frac{\sqrt{P} \cdot \sin L}{1 + E\cos V}
\]
oder
\[
\frac{\sqrt{p} \cdot \sin l}{r} = \frac{\sqrt{P} \cdot \sin L}{R}
\]
enthalten.

\newpage

Der Beweis st\"utzt sich auf folgende Figur:

$A$ Ort der Erde, $B$ geoc., $C$ helioc.~Ort
des Planeten; $A'$, $C'$ Pl\"atze von $A$ und $C$ nach
unendl.~kl.~Zwischenzeit $dt$.

Es ist also
\[
AA' = R\,dL, \quad CC' = r\,dl.
\]

Verm\"oge der Natur des Problems m\"ussen die beiden neuen Pl\"atze von der durch
Sonne, Erde und Planet gelegten [durch $ACB$ repr\"as.~Ebene] gleich weit ab-
stehen. Diese Abst\"ande sind $R\,dL\sin A$ und $r\,dl\sin C$; da nun $\sin A : \sin C
= \sin l : \sin L$, so wird $\frac{R\,dL}{r\,dl} = \frac{\sin L}{\sin l}$, oder da $RR\,dL : rrl\,dl = \sqrt{P} : \sqrt{p}$:
\[
\frac{\sqrt{p} \cdot \sin l}{r} = \frac{\sqrt{P} \cdot \sin L}{R},
\]
welches die obige Gleichung ist.

\subsection*{[10. Zu Art. 140.]}

I. Es seien $\mathfrak{A}$, $\mathfrak{A}'$, $\mathfrak{A}''$ die Abst\"ande der drei Punkte $A$, $A'$, $A''$ von dem
gr\"ossten Kreise $BB'B''$; $h$ die Neigung dieses gr\"ossten Kreises gegen die
Ekliptik, und $H$ die L\"ange seines aufsteigenden Knotens auf der Ekliptik.
Man hat dann
\begin{align*}
\sin\mathfrak{A} &= \sin h \sin(l - H) \\
\sin\mathfrak{A}' &= \sin h \sin(l' - H) \\
\sin\mathfrak{A}'' &= \sin h \sin(l'' - H);
\end{align*}
woraus leicht mit H\"ulfe des Lemma I in Art.~78 folgt:
\begin{equation}\tag{1}
\sin\mathfrak{A}\sin(l' - l'') + \sin\mathfrak{A}'\sin(l'' - l) + \sin\mathfrak{A}''\sin(l - l') = 0.
\end{equation}

II. Es seien ferner $B$, $B'$, $B''$ die Winkel, welche der gr\"osste Kreis
$BB'B''$ mit $AB$, $A'B'$, $A''B''$ macht, wodurch
\begin{align*}
\sin\mathfrak{A} &= \sin\delta \sin B \\
\sin\mathfrak{A}' &= \sin(\delta' - \alpha')\sin B' \\
\sin\mathfrak{A}'' &= \sin\delta''\sin B''.
\end{align*}

\newpage

Die Gleichung 1 erh\"alt demnach folgende Gestalt:
\begin{equation}\tag{2}
\sin\delta\sin B\sin(l' - l'') + \sin(\delta' - \alpha')\sin B'\sin(l - l'') + \sin\delta''\sin B''\sin(l'' - l) = 0.
\end{equation}

III. Man hat ferner
\[
\frac{AD' - \delta}{\sin B} = \frac{BD'}{\sin\delta}, \quad \frac{A''D' - \delta''}{\sin B''} = \frac{B''D'}{\sin\delta''},
\]
folglich
\begin{equation}\tag{3}
\sin\delta\sin B = \frac{BD'}{AD' - \delta}, \quad \sin\delta''\sin B'' = \frac{B''D'}{A''D' - \delta''}
\end{equation}
wie auch
\[
\frac{A''D - \delta}{\sin B''} = \frac{B''D}{\sin\delta''}, \quad \frac{A'D' - \delta'}{\sin B'} = \frac{B'D'}{\sin\delta'},
\]
folglich
\begin{equation}\tag{4}
\sin B' = \frac{B'D'}{A'D' - \delta'}, \quad \sin B'' = \frac{B''D}{A''D - \delta}
\end{equation}
und
\[
\frac{B''\sin\delta''}{B'\sin\delta'} = \frac{A''D' - \delta''}{A'D' - \delta'}.
\]

Diese Werthe aus 3, 4 in 2 substituirt geben eine Gleichung, die bei
n\"aherer Betrachtung mit der ersten Gleichung am Schluss von Art.~140 iden-
tisch gefunden wird.

IV. Um zum Beweise der andern Formel zu gelangen, sei $P$ der Pol
der Ekliptik, $\Pi$ der Punkt der Ekliptik, dessen L\"ange oben mit demselben
Buchstaben bezeichnet wurde, also mit $BB''$ in Einem gr\"ossten Kreise,
also
\[
HP = 90^\circ, \quad BP = 90^\circ, \quad B''P = 90^\circ, \quad BHP = 90^\circ, \quad BPB'' = 90^\circ - \alpha'' + \alpha.
\]

Im Dreiecke $BPB''$ [dreieck $HB''P$]:
\[
\sin HB \cdot P = \cos B \cdot \sin(\alpha'' - \alpha)
\]
folglich
\begin{equation}\tag{5}
\sin(\alpha'' - \alpha) = \frac{\sin HB \cdot P}{\cos B \cos B''}.
\end{equation}

Nun ist nach Art.~138, [$\delta'$]
\[
\delta = \tan\beta \sin(\alpha'' - l') - \tan\beta''\sin(\alpha - l').
\]

\newpage

also, da leicht zu \"ubersehen ist, dass
\[
\tan\beta = \tan h \sin(\alpha - H), \quad \tan\beta'' = \tan h \sin(\alpha'' - H):
\]
\[
\delta = \tan h \left\{\sin(\alpha'' - l')\sin(\alpha - H) - \sin(\alpha - l')\sin(\alpha'' - H)\right\}
\]
oder, da nach Lemma I Art.~78 (wenn dort statt $A$, $B$, $C$ \ldots\ $\alpha - l'$, $\alpha'' - l'$,
$H - l'$ gesetzt wird) identisch ist
\[
\sin(\alpha - H)\sin(\alpha'' - l') + \sin(l' - \alpha'')\sin(\alpha - l') + \sin(\alpha'' - \alpha)\sin(H - l') = 0:
\]
\[
\delta = \tan h \cdot \sin(\alpha'' - \alpha)\sin(l' - H),
\]
oder, f\"ur $\sin(\alpha'' - \alpha)$ den Werth aus 5 substituirt,
\[
\delta = \frac{\sin h \sin(l' - H) \sin BB''}{\cos B \cos B''}.
\]

V. Wir haben ferner
\[
\sin\mathfrak{A}' = \sin h \sin(l' - H) = \sin(\delta' - \alpha')\sin B'.
\]
Es ist also (Schluss von Art.~140)
\[
\sin B : \sin B'' = \sin(\alpha'' - \alpha) : \sin(l' - H)
\]
oder da
\[
\delta = \frac{\sin h \sin(l' - H) \sin BB''}{\cos B \cos B''} = \frac{\sin\mathfrak{A}' \sin BB''}{\cos B \cos B''}
\]
d.~i.
\[
\delta = \frac{\sin\mathfrak{A}'}{\cos B \cos B''} \cdot \sin BB''
\]
oder
\[
V = \frac{\sin\delta' \sin BB''}{\cos B \cos B'' \sin B''} = \frac{\sin\delta' \sin(AD' - \delta)}{\sin B'' \cos B \cos B''}
\]
oder
\[
V = \frac{\sin\delta' \sin(AD' - \delta) \sin e'}{\sin B' \cos B \cos B'' \sin(AD' - \delta')}
\]
identisch mit der Gleichung am Schluss von Art.~140.

\subsection*{[11. Zu Art. 182--184.]}

\textit{Zeitschrift f\"ur Astronomie, herausgegeben von B.~von Lindenau und J.~G.~F.}

\textit{Erster Band. Heft f\"ur M\"arz und April 1816.}

Ich zeige bei dieser Gelegenheit noch eine Berichtigung an, die S.~218
und 219 [S.~250--252 dieses Bandes] der \textit{Theoria Motus Corporum Coelestium}
zu machen ist. Man lese nemlich

S.~218 Z.~3 statt $e \ldots\ e''' \ldots\ e'''''$

Z.~4 statt $\sqrt{A}, \sqrt{B'}, \sqrt{C''}, \sqrt{D'''}$ der Ordnung nach $\sqrt{A'}, \sqrt{B''}, \sqrt{C'''}, \sqrt{D^{IV}}$.

S.~219 Z.~8 statt $\sqrt{A}, \sqrt{B'}, \sqrt{C''}, \sqrt{D'''}$ der Ordnung nach $\sqrt{A'}, \sqrt{B''}, \sqrt{C'''}, \sqrt{D^{IV}}$.

Diese Unrichtigkeit hat ihren Ursprung in dem Umstande, dass fr\"uher
bei der Untersuchung andere Bezeichnungen angewandt waren, deren Umtausch
an den angezeigten Stellen vers\"aumt war. Das numerische Beispiel ist, wie
man sieht, den berichtigten Ausdr\"ucken gem\"ass berechnet, doch ist darin ein
Rechnungsfehler begangen. Die Gleichung S.~219 Z.~5 v.~u.~muss nemlich sein
\[
6633 = 12707 + 2P - 9Q + 123R,
\]
folglich die Genauigkeit von $r$

\newpage

\begin{center}
S.~220
\end{center}

\begin{center}
\begin{tabular}{c|c|c}
$\sqrt{2P}$ & $\sqrt{2Q}$ & $\sqrt{2R}$ \\
\hline
$211$ & $784$ & $156$
\end{tabular}
\end{center}

Ich bin auf diese Berichtigungen durch Herrn Nicolai aufmerksam ge-
macht worden, welchem ich daf\"ur hier den verbindlichsten Dank sage.

\subsection*{[12. Zu Art. 183.]}
\textbf{Bestimmung der wahrscheinlichen Fehler der Resultate
der Methode der kleinsten Quadrate.}

Man setze, der Algorithmus, welcher im I.~Bde der G[\"ottinger] C[ommenta-
tiones]: \textit{Disquisitio de elementis ellipticis Palladis}, [Werke, Band VI, S.~21--22]
gelehrt ist, f\"uhre auf folgende Gleichungen zur Bestimmung der unbekannten

\begin{center}
$39^*$
\end{center}

Gr\"ossen $p$, $q$, $r$, $s$
\begin{align*}
[A =]\quad 0 &= X + ap + \beta q + \gamma r + \delta s \\
[B =]\quad 0 &= X' + \beta'q + \gamma'r + \delta's \\
[C =]\quad 0 &= X'' + \gamma''r + \delta''s \\
[D =]\quad 0 &= X''' + \delta'''s.
\end{align*}

Man bestimme $\alpha$, $\beta$, $\gamma$, $\delta$ durch die Gleichungen
\begin{align*}
\alpha^2 &= \lambda^2 \text{ (Wahrsch.~Beob.~Fehler)} \\
\tau + \tau'x + \tau''p &= u \\
\sigma + \sigma'q + \sigma''p + \sigma'''s &= 0,
\end{align*}
so ist der wahrscheinliche Fehler von $p$:
\[
\sqrt{\frac{n}{n - 4}(\alpha^2 + \tau'xx + \tau''pp + \sigma'''ss)}.
\]

Noch zierlicher so: Aus den Gleichungen
\begin{align*}
\alpha p + \beta q + \gamma r + \delta s &= u \\
\beta'q + \gamma'r + \sigma's &= u' \\
\gamma''r + \sigma''s &= u'' \\
\sigma'''s &= u'''
\end{align*}
entstehe durch Elimination
\begin{align*}
p &= p'u + p''u' + p'''u'' + p^{IV}u''' \\
q &= q'u' + q''u'' + q'''u''' \\
r &= r''u'' + r'''u''' \\
s &= s'''u''';
\end{align*}
sodann sind die wahrscheinlichsten Werthe
\begin{align*}
-P &= p'X + p''X' + p'''X'' + p^{IV}X''' \\
-q &= q'X' + q''X'' + q'''X''' \\
-r &= r''X'' + r'''X'''
\end{align*}
und die wahrscheinlichen Fehler dieser Bestimmungen, den wahrscheinlichen
Fehler der Beobachtungen $= 1$ gesetzt:

\newpage

\begin{center}
\begin{tabular}{ll}
f\"ur $p$ & $\sqrt{(p'p' + p''p'' + p'''p''' + p^{IV}p^{IV})}$ \\
f\"ur $q$ & $\sqrt{(q'q' + q''q'' + q'''q''')}$ \\
f\"ur $r$ & $\sqrt{(r''r'' + r'''r''')}$ \\
f\"ur $s$ & $s'''\sqrt{(s'''s''')}$.
\end{tabular}
\end{center}

\subsection*{[13.]}
\textbf{Zusammenhang des Initialzustandes mit den Elementen [einer Planetenbahn].}

\begin{center}
\begin{tabular}{ll}
$r$ & Radius Vector \\
$V$ & L\"ange in der Bahn \\
$\lambda$ & L\"ange in der Ekliptik \\
$\beta$ & Breite \\
$\Omega$ & L\"ange des aufsteigenden Knotens \\
$i$ & Neigung der Bahn \\
$e = \sin\varphi$ & Excentricit\"at \\
$p$ & halber Parameter \\
$\omega$ & Perihelium \\
$t$ & Zeit \\
$k$ & Constannte $= \sqrt{M}$
\end{tabular}
\end{center}

\begin{align}
\tag{1} \tan\beta &= \tan i \sin(\lambda - \Omega) \\
\tag{2} \frac{r\cos V\,dV}{k\,dt} &= \sqrt{p} \\
\tag{3} \frac{r\,d\lambda}{k\,dt} &= \sqrt{p} \cdot \cos i \\
\tag{4} \frac{rr\,du}{k\,dt} &= \sqrt{p} \\
\tag{5} \frac{2}{r} - \frac{du^2 + d\beta^2}{dt^2} \cdot \frac{1}{k^2} &= \frac{1}{a} \\
\tag{6} \frac{rr\,du}{k\,dt} &= \sqrt{p} \\
\tag{7} 2k \arctan\sqrt{\frac{1 - e}{1 + e}} \tan\tfrac{1}{2}(V - \omega) &= \int \frac{k\,dt}{rr} = u - \varepsilon \\
\tag{8} \frac{1}{r} &= \frac{1 + e\cos(V - \omega)}{p} = \frac{1}{a}
\end{align}

\newpage

\subsection*{[14.]}
\textbf{\"Uber die Verbesserung der Elemente eines Planeten
durch vier beobachtete Oppositionen.}

Wenn die Elemente eines Planeten schon so genau bekannt sind, dass die angenommenen Werthe
von den wahren nur sehr wenig abweichen und man hat die Beobachtungen von vier Oppositionen, so l\"asst
sich die Verbesserung der Elemente leichter und einfacher vollbringen, als wenn die vier gegebenen \"Orter
des Planeten beliebig genommen w\"aren. Das Theoretische dieser Untersuchung ist schon in der Abhandlung
\textit{``Disquisitio de elementis ellipticis Palladis''} mitgetheilt; hier folgen einige praktische Bemerkungen, die \"uber-
haupt bei der Berechnung der Oppositionen und dergl.~brauchbar sind.

Zuerst muss die Epoche verbessert werden. Man berechnet also aus einer beliebigen Beobachtung,
am besten aus einer solchen, die zu einer der mittlern Oppositionen geh\"ort, die geocentrische L\"ange und
Breite des Planeten, aus diesen mit H\"ulfe des aus der unver\"anderten Epoche durch die bekannten Formeln
erhaltenen Radius Vector, die heliocentrische L\"ange und Breite, aus diesen die L\"ange in der Bahn. Nach-
dem an diese L\"ange die St\"orungen wenn schon Formeln f\"ur sie vorhanden sind; angebracht, leitet man
aus ihr die mittlere Anomalie her, welche mit der aus der unver\"anderten Epoche erhaltenen verglichen, die
\"Anderung der Epoche gibt. Sind keine Beobachtungen, sondern nur die berechneten L\"angen und Breiten
zu der Zeit der vier Oppositionen gegeben, so wird man eine dieser L\"angen zu w\"ahlen haben; besser ist
es indessen, die Beobachtungen selbst vorzunehmen, als sich auf die Berechnungen Anderer zu verlassen.

Mit der auf solche Art verbesserten Epoche und den \"ubrigen unver\"andert gelassenen Elementen
werden nun die vier Oppositionen vorl\"aufig berechnet, wobei man sich mit hinl\"anglicher Genauigkeit fr\"u-
erer Berechnungen bedienen kann. Um eine Opposition am bequemsten berechnen zu k\"onnen, verfahrt
man folgendermaassen.

Man w\"ahlt, wenn von mehrem Orten Beobachtungen da sind, die Beobachtungen desjenigen Orts,
der die meisten geliefert hat, bestimmt aus den angegebenen Zeiten des Durchgangs durch den Mittagskreis,
die Culminationszeiten des Planeten von zwei zu zwei Tagen f\"ur die ganze Zeit, welche alle Beobachtungen
umfassen. Durch eine vorl\"aufige Rechnung, f\"ur einen in der Mitte der Beobachtungszeit liegenden Tag,
bestimmt man die Entfernung des Planeten von der Erde und aus dieser die Zeit, welche das Licht ge-
braucht, um von dem Planeten auf die Erde zu kommen, welche Zeit von den Culminationszeiten abzuziehen
ist; hiedurch hat man f\"ur die Aberration Rechnung gehalten. Aus der verbesserten Epoche leitet man
nun durch Abziehung der f\"ur die Zeit der Epoche stattfindenden L\"ange des Perihels die mittlere Anomalie
her und bestimmt aus dieser mittelst der t\"aglichen tropischen Bewegung die mittlern Anomalien f\"ur die an-
genommenen Tage und aus diesen die Rectascensionen und Declinationen. Um f\"ur die Nutation Rechnung
zu halten, wendet man nicht die mittlere durch die Elemente gegebene L\"ange des Knotens, sondern die
wahre an, die man durch Hinzuf\"ugung der Nutation in der L\"ange Bohnenbergers [Astron.~p.~672] zu jener
erh\"alt; die Parallaxe aber bringt man nach den in der \textit{Theoria motus Corpor.~Coel.}~art.~59 oder 70 er-
kl\"arten Methode an. \"Ubrigens sieht man leicht ein, weshalb die Rechnung von zwei zu zwei Tagen durch-
gef\"uhrt ist, auch wenn an einigen derselben nicht sollte beobachtet sein; man kann n\"amlich die Rechnung
sehr bequem durch Differenzen pr\"ufen, so dass es nicht m\"oglich ist, einen Fehler zu \"ubersehen, und dann
vermittelst einer scharfen Interpolation die Berechnungen f\"ur alle Tage erhalten.

Die auf solche Weise erhaltenen Rectascensionen und Declinationen vergleicht man mit den beob-
achteten, bestimmt hieraus die mittlere Abweichung sowohl der Rectascension als der Declination und
indem man diese zu den berechneten hinzuf\"ugt, leitet man f\"ur den Tag der Opposition, welcher sehr bald

\newpage

durch Vergleichung der Beobachtungen mit den berechneten Sonnenl\"angen gefunden wird, die geocentrische
L\"ange her, bestimmt eben dieselbe f\"ur einen um eine beliebige Gr\"osze (etwa einige Stunden) von der vorhin
angenommenen Zeit abweichenden Zeitpunkt und sucht nun durch Interpolation den Augenblick zu be-
stimmen, in welchem die L\"ange der Sonne und die des Planeten genau um $180^\circ$ von einander verschieden
sind. Bei der Berechnung der Sonne ist \"ubrigens f\"ur die durch die Aberration des Planeten verbesserte
Culminationszeit die wahre L\"ange zu suchen, also die durch die Tafeln gefundene, welche die Aberration
der Sonne bereits einschliesst, um diese zu vermehren.

Sind nur wenige Beobachtungen vorhanden (etwa bis sechs), so kann man auch, wenn man es f\"ur
vortheilhafter h\"alt, die Rectascensionen und Declinationen bloss f\"ur die Beobachtungstage berechnen.

Bei der vorl\"aufigen Berechnung der Oppositionen kann man allenfalls die Aberration, Nutation und
Parallaxe vernachl\"assigen; will man aber die Oppositionen nur einmal berechnen, so muss nothwendig auf
sie R\"ucksicht genommen werden.

Aus den gefundenen L\"angen und Breiten zu der Zeit der vier Oppositionen leitet man nun nach der
in der Abhandlung \textit{de Elementis Palladis} angef\"uhrten Methode die Verbesserungen der einzelnen Ele-
mente her.

\subsection*{[15.]}
\textbf{Differentialformeln bei Berechnung der Oppositionen.}

\begin{center}
\begin{tabular}{ll}
$\lambda$ & ber.~hel.~L\"ange \\
$\lambda'$ & beob.~hel.~L\"ange \\
$\omega$ & Sonnenn\"ahe \\
$V$ & wahre Anomalie \\
$r$ & Radius Vector \\
$a$ & halbe grosse Axe \\
$i$ & Neigung der Bahn \\
$\Omega$ & Knoten \\
$n$ & Tage seit Epoche \\
$\gamma$ & t\"agliche Bewegung \\
$E$ & excentrische Anomalie \\
$\delta$ & heliocentrische Breite \\
$e = \sin\varphi$ & Excentricit\"at \\
$\beta$ & geocentrische Breite
\end{tabular}
\end{center}

\begin{align*}
\sin 2(\lambda - \Omega) &= A \\
\sin 2(V + \omega - \Omega) &= B \\
\frac{1 - e\cos E}{r} &= C
\end{align*}

\begin{align*}
d\lambda &= E\,d\text{Ep.} + (nE - \tfrac{\delta}{m})\,d\tau + (B - E)\,d\omega + F\,d\varphi + C\,di + (D - E)\,d\Omega - \tan\delta \sin 2(\lambda - \Omega)\,di \\
d\delta &= E'\,d\text{Ep.} + (nE' - \tfrac{\varepsilon}{m})\,d\tau + (B' - E')\,d\omega + F'\,d\varphi + C'\,di + (D' - E')\,d\Omega + \ldots
\end{align*}

\newpage

\textbf{Beispiel. Opposition der Pallas von 1807.}

[Mit den in der \textit{Disquisitio de elementis ellipticis Palladis} art.~2, 3, 8 ge-
gebenen Werthen
\begin{center}
\begin{tabular}{ll}
$\log a = 0{,}41223$ & \\
$\psi = 107^\circ 47' 31''$ & \\
$e = 0{,}24176$ & \\
$\lambda = 223^\circ 37' 25''$ & \\
$\delta = 1^\circ 4' 0{,}4''$ & \\
$\omega = 223^\circ 3' 7' 28''$ & \\
$\beta = 42^\circ 1' 1' 28''$ & \\
$\Omega = 172^\circ 3' 2' 24''$ & \\
$i = 34^\circ 3' 7' 321''$ & \\
$\delta = 28^\circ 14' 51''$ &
\end{tabular}
\end{center}
[und den leicht zu findenden
\[
E = 93^\circ 4' 6' 4'', \quad \log r = 0{,}44916
\]
ergibt die Rechnung:]

\begin{center}
\begin{tabular}{lll}
$\sin 2(\lambda - \Omega) \ldots 9{,}99013$ & $\sin\beta \ldots 9{,}82712$ & $\sin\beta \ldots 9{,}82712$ \\
$\sin 2(V + \omega - \Omega) \ldots 9{,}96466$ & $\sin(\beta - \delta) \ldots 9{,}38$ & 
\end{tabular}
\end{center}

\newpage

\subsection*{[16.]}
\textbf{Grenzen der geocentrischen \"Orter.}

\textit{(Band VI, S.~106 ff.)}

Die Theorie der Grenzen der geocentrischen \"Orter theilt die Oberfl\"ache
der Himmelskugel in verschiedene Theile, und die nach der von mir gegebenen
Methode bestimmten Linien sind, allgemein zu reden, nichts anderes, als Scheidungen
zwischen zwei Fl\"achentheilen der Kugel, wo die auf der einen Seite liegenden
Punkte auf zwei Arten mehr geocentrische \"Orter sein k\"onnen als die auf der

\newpage

andern. Die in den Scheidungslinien liegenden Punkte machen den \"Uber-
gang, d.~i.~sie sind auf Eine Art mehr als die Punkte auf der einen Seite,
und auf Eine Art weniger als die Punkte auf der andern Seite f\"ahig geocen-
trischer \"Orter zu sein. \"Ubrigens l\"asst sich auch das Criterium angeben, wo-
nach a priori entschieden wird, auf welcher Seite der Scheidungslinie zwei
Aufl\"osungen mehr stattfinden als auf der andern, wobei ich jedoch gegen-
w\"artig mich nicht aufhalten will[$^*$].

Alle bisher bekannten Planeten haben solche Bahnen, dass die durch die
Theorie gefundenen Scheidungslinien immer wahre Limiten sind. Es sind
nemlich zwei Scheidungslinien, die die Himmelskugel in drei Fl\"achenr\"aume
abtheilen; zwei sind isolirte Fl\"achen, in welche gar keine geocentrische \"Orter
fallen, der dritte zwischen jenen, g\"urtelartig, enth\"alt alle geocentrischen \"Orter
und jeder Punkt innerhalb des G\"urtels kann auf zwei Arten, jeder Punkt
auf der Limite auf Eine Art geocentrischer Ort sein.

Es lassen sich aber fingirte Bahnen denken [$^{**}$)] (Cometenbahnen werden
vielleicht mehrere in dem Fall sein, wor\"uber ich jedoch bisher keine Nach-
forschung gemacht habe), wo zwei Limiten die Himmelskugel auch in drei
St\"ucke scheiden, und wo der eine Theil gar keine geocentrische \"Orter ent-
h\"alt, der andere die geocentrischen \"Orter zu zwei Arten, der dritte hingegen
die Punkte, die auf vier verschiedene Arten geocentrische \"Orter sein k\"onnen.

Noch weitere Mannigfaltigkeit ergibt sich, wenn Eine Scheidungslinie sich
selbst einmal oder zweimal schneidet, eine einfache oder doppelte Schleife
bildet. Im ersten Fall wird sie zwei, im zweiten drei Fl\"achenr\"aume von dem
g\"urtelf\"ormigen Theile abtrennen, in denen resp.~4 und 0 oder 4, 0, 4 Auf-
l\"osungsarten stattfinden.

Noch anders verh\"alt es sich mit einer Bahn, die in die Erdbahn wie ein
Kettenring eingreift. In einem solchen Fall ist nur Eine zusammenh\"angende
Limitenlinie, also zwei geschiedene Fl\"achentheile, wo die Punkte des einen
Theils auf Eine Art, die des andern Theils auf drei Arten geocentrische \"Orter
sein k\"onnen.

\begin{flushright}
[$^*$) Siehe Notiz 17.]

[$^{**}$) Siehe Notiz 18.]
\end{flushright}

\newpage

\subsection*{[17.]}
\textbf{Zodiakus auf der Himmelskugel.}

\textit{1847 Oct.~28.}

F\"ur jeden Punkt der Himmelskugel ist entweder gar keine, eine, oder
mehr als eine Combination vorhanden, die die Sichtbarkeit daselbst des einen
Planeten vom andern aus bedingt, oder eine bestimmte Zahl von Aufl\"osungen.
Beim Durchgange durch die Limitenlinie \"andert sich allemal die Anzahl der
Aufl\"osungen um 2 Einheiten; so lange aber, bei irgend einem Wege auf der
Himmelskugel, die Limitenlinie nicht getroffen wird, bleibt die Zahl der Auf-
l\"osungen unge\"andert. Die Frage ist nun, wie man die Seite der Limite, wo
zwei Aufl\"osungen mehr sind, von der andern, wo zwei weniger sind, unter-
scheidet. Folgende Betrachtungen dienen dazu.

I. Es sei $p$ ein Punkt der Planetenbahn, von welcher aus die Punkte
der andern Bahn auf die Himmelskugel projicirt werden; $P$ der correspon-
dirende Punkt der andern Bahn; $\Pi$ der Punkt der Knotenlinie, wo die Tan-
genten an den beiden Planetenbahnen in $p$ und $P$ einander schneiden;
ferner sei $S$ ein anderer beliebiger fester Punkt der Knotenlinie (etwa die Sonne).
Die drei Figuren hat man sich in drei verschiedenen Ebenen vorzustellen: I~in
der Ebene durch $\Pi pP$; II~in der Ebene durch $\Pi Sp$; III~in der Ebene
durch $\Pi SP$. Noch bezeichnen wir mit

$p$, $P$ die Winkel der Ebenen II und III mit I,

$q$, $Q$ die Winkel $p\Pi S$, $P\Pi S$,

$u$, $U$ die Winkel $\Pi pP$, $\Pi Pp$,

$r$, $R$ die Kr\"ummungshalbmesser der beiden Bahnen in $p$ und $P$.

Es seien nun $p'$, $P'$ zwei andere Punkte correspondirend, und dazu der
Punkt $\Pi'$ der Knotenlinie. $\Pi\Pi' = \zeta$ ist als unendlich klein angenommen. Es
werden dann die Normalen aus $p'$, $P'$ auf $pP$:
\[
p'\pi' = n\zeta, \quad P'\Pi' = N\zeta.
\]

\newpage

Ferner seien $p''$, $P''$ zwei andere Punkte in den Tangenten $p\Pi$, $P\Pi$, so dass
$p''P''$ parallel mit $pP$, und $\frac{\zeta}{\omega} = \frac{p\Pi}{p''p} = \frac{P\Pi}{P''P} = \omega$ unendlich klein. Eine durch $p''P''$
normal gegen I gelegte Ebene treffe die Curven in $p'''$, $P'''$, und $p^\text{iv}$, $P^\text{iv}$ seien
die Fusspunkte der von $p'''$, $P'''$ auf I gef\"allten Perpendikel. Man hat dann
\[
p'''p^\text{iv} : P'''P^\text{iv} = p'P' : pP \cdot \frac{n\sin u}{N\sin U}
\]
\[
\frac{p''p^\text{iv}}{p'''p^\text{iv}} = \frac{P''P^\text{iv}}{P'''P^\text{iv}} \cdot \frac{pP}{p'P'} \cdot \frac{\sin p}{\sin P} \cdot \frac{\sin q}{\sin Q} \cdot \frac{\sin u}{2r}
\]
\[
\frac{p'''p^\text{iv} \cdot p''p^\text{iv}}{P'''P^\text{iv} \cdot P''P^\text{iv}} = \frac{n\sin u \cdot \sin p \sin q \sin u}{N\sin U \cdot \sin P \sin Q \sin U}.
\]

Man sieht aber leicht, dass $n\sin u = N\sin U$ das Perpendikel aus $\Pi$ auf
$pP$; ferner $\sin p \sin q = \sin P \sin Q = \text{Sinus des Winkels, welchen die Linie}
\Pi S \text{ mit der Ebene I, und folglich mit dem gedachten Perpendikel macht,}$
woraus man leicht schliesst, dass $n\sin p \sin q \sin u = N\sin P \sin Q \sin U$ nichts
anderes ist, als die k\"urzeste Entfernung der Geraden $pP$, $\Pi S$ von einander.

\newpage

Endresultat ist also
\[
\frac{p'''p^\text{iv} \cdot p''p^\text{iv}}{P'''P^\text{iv} \cdot P''P^\text{iv}} = \frac{p'P'}{pP}.
\]
Ist folglich $p'''p^\text{iv} > P'''P^\text{iv}$, so ist $p'p < P'P$ und umgekehrt. Die erste
Ungleichheit zeigt aber an, dass eine mit $p'''P'''$ parallel aus $p$ gezogene Ge-
rade einerseits und $S$ andererseits auf entgegengesetzte Seite von I fallen;
die zweite, dass die Projectionen auf die Himmelskugel von den beiden Ge-
raden $p'P'$ und $pP$, wovon die letztere nur um ein unendlich kleines $2^\text{ten}$ Ord-
nung von der Limitenlinie absteht auf Einer Seite von der Projection von
$pP$ liegen.

Hieraus ist das gew\"unschte Criterium leicht abzuleiten; man kann es so
vortragen:

[Sei] $B$ [die] \"aussere Planetenbahn von der Sonne aus auf die Himmels-
kugel bezogen; dieser gr\"osste Kreis theilt die Himmelskugel in zwei H\"alften
$\mathfrak{A}$ und $\mathfrak{B}$. Es sei $\mathfrak{A}$ diejenige, in welcher der heliocentrische Ort des In-
ern Planeten erscheint. Dann liegt gegen ein St\"uck der Limite $LL'$ diejenige
Gegend, wo zwei Aufl\"osungen mehr gelten, eben so wie $\mathfrak{A}$ gegen $B$, indem
man in $LL'$ und $B$ als vorw\"arts gehenden Sinn Zusammengeh\"origes annimmt.
Es ist hiebei angenommen, dass der Planet der \"aussern Bahn der beobachtete
ist. Im entgegengesetzten Fall ist es umgekehrt.

Symbolisch und noch allgemeiner kann das Criterium so ausgedr\"uckt
werden:

Es seien $l$, $l'$ zwei einander unendlich nahe Punkte in einer Limite; $p$, $p'$
die entsprechenden heliocentrischen \"Orter des Planeten, von welchem aus [be-
obachtet wird], $P$, $P'$ hingegen die des Planeten, welcher beobachtet wird; $a$
ein Punkt nahe an $l'$ auf der Seite, wo zwei Aufl\"osungen mehr gelten als
auf der andern. Dann ist, wenn das Congruenzzeichen $\equiv$ zur Bezeichnung
gleicher Lageordnung dient, und ein vorgesetztes $-$~Zeichen die entgegen-
gesetzte Lageordnung bedeutet:
\[
\Pi'a \equiv PPp \equiv -pp'P.
\]

\subsection*{[18.]}
\textbf{Zodiak[u]s.}

[Um ein Beispiel f\"ur den Fall zu erhalten, wo zwei Limiten die Himmels-
kugel in mehrere St\"ucke scheiden, und wo einige dieser St\"ucke gar keine
geocentrische \"Orter enthalten, andere solche, die auf zwei Arten, wieder
andere solche, die auf vier Arten geocentrische \"Orter sein k\"onnen, w\"ahle man
f\"ur die]

\textbf{\"Aussere Bahn} $k = 1$, $e = 0$

\textbf{Innere Bahn} $k' = 0{,}8$, $\varphi' = 110^\circ$

[wo die Bezeichnungen dieselben sind, wie in der Abhandlung >>\"Uber die
Grenzen der geocentrischen \"Orter etc.<< Band VI S.~106. Man hat dann die]

\textbf{Formeln:}
\begin{align*}
r'\cos t' &= k\cos t \\
r'\sin t' &= k\sin t \\
r'\cos t' - k\cos t &= A\cos\alpha\cos\beta \\
-k\sin t &= A\sin\alpha\cos\beta \\
r'\sin t' &= A\sin\beta.
\end{align*}

[In diesem Falle existiren zwei getrennte Limitenlinien, welche sich selbst
schneiden k\"onnen, und zwar findet sich die] kritische Stelle [f\"ur]

\newpage

\subsection*{[19.--20.]}
\textbf{Zur parabolischen Bewegung.}

\subsection*{[1.]}
[Nach den Bezeichnungen der \textit{Theoria motus} Art.~98 und 108 hat man:
\[
\xi = \frac{r' + r + p}{2\sqrt{rr'}} = 1 + \cos f\cos F, \quad \eta = \frac{r' - r}{2\sqrt{rr'}} = \sin f\sin F,
\]
also
\[
\xi\eta = (\cos f + \cos F)\sin f\sin F
\]
oder, da
\[
\frac{p}{rr'} = \sin f\sin F,
\]
wo
\[
\xi = \frac{1 + \cos f^2 + 2\cos f\cos F}{2(\cos f + \cos F)}, \quad \eta = \frac{\sin f^2\sin F}{2(\cos f + \cos F)},
\]
so wird
\[
\frac{p}{r' + r} = \frac{\sin f\sqrt{(1 + \cos f^2 + 2\cos f\cos F)}}{1 + \cos f\cos F}, \quad \frac{r' - r}{p} = \frac{\sin F}{\sqrt{(1 + \cos f^2 + 2\cos f\cos F)}}.
\]

Hieraus:
\[
\cos f + \cos F = \frac{\cos f(1 + \cos f^2 + 2\cos f\cos F)}{1 + \cos f\cos F},
\]
\[
\cos\psi\tan\varphi = \tan f,
\]
woraus $f$ gefunden wird.

Ferner
\[
\sin\psi\sin f = \frac{\sin f(\cos f + \cos F)}{1 + \cos f\cos F},
\]
\[
\sin\psi\sin\vartheta = \frac{\sin f\sin F}{1 + \cos f\cos F},
\]
hieraus
\begin{align*}
\sin\psi\sin(F - \psi) &= \cos f\sin\psi\sin\varphi = \sin f\tan\psi\cos\psi \\
\sin\psi\cos(F - \psi) &= \sin f\sin\psi.
\end{align*}

Diese Gleichungen selbst, oder die daraus folgenden:
\[
\tan\psi\cos\psi = \tan(F - \psi)
\]
oder
\begin{align*}
\sin f\sin\psi\tan\tfrac{1}{2}\varphi &= \sin(2\psi - F) \\
\sin f\sin\psi\cot\tfrac{1}{2}\varphi &= \sin F
\end{align*}
[dienen zur Bestimmung von $F$;] die wahren Anomalien [sind dann] $F - f$
und $F + f$.

[Weiter ist
\[
\frac{r' + r}{2p} = \frac{\sin f}{\sin\psi\sin\varphi\cos\tfrac{1}{2}\varphi} = \frac{\tan\tfrac{1}{2}\varphi}{\sin\psi\cos\tfrac{1}{2}\varphi\cos\tfrac{1}{2}\psi},
\]
also]
\[
q = \frac{p}{1 + \cos\varphi} = \tfrac{1}{2}(r' + r)\cos\tfrac{1}{2}\varphi \cdot \cos\tfrac{1}{2}\psi,
\]
woraus $q$ bekannt ist.

Zur Berechnung der seit dem Periheldurchgange verflossenen Zeiten hat
man mit der Bezeichnung der Notiz 1:
\begin{align*}
mm &= 2(r' + r)\cos(45^\circ - \tfrac{1}{2}\varphi)^2 \\
nn &= 2(r' + r)\sin(45^\circ - \tfrac{1}{2}\varphi)^2, \quad \text{also} \quad m + n = 2\sqrt{(r' + r)}\cdot\cos\tfrac{1}{2}\varphi \\
mn &= (r' + r)\cos\varphi \quad m - n = 2\sqrt{(r' + r)}\cdot\sin\tfrac{1}{2}\varphi.
\end{align*}
Da nun
\[
r' - r = (r' + r)\sin\psi\sin\varphi,
\]
so wird]

\newpage

\begin{equation}\tag{4}
\frac{t' - t}{k} = \frac{mmn - nnn}{12kk'} \cdot \frac{4(r' - r) + 3(m - n)^2(mn + nn) - (m - n)^2(mm + mn + nn)}{m}
\end{equation}
Letztere Formel allein findet sich leichter so:
\[
t' - t = \frac{1}{6kk'}(r' + r + p)^{\frac{3}{2}} - (r' + r - p)^{\frac{3}{2}} = \frac{1}{6kk'}(r' + r)^{\frac{3}{2}} \cdot \ldots
\]

\subsection*{[2.]}
\textbf{Zur Berechnung der Chorde $p$ aus der Zeit $= T$ [$= t' - t$], und der Summe
der beiden Radien $= r + r'$.}

\[
6kT = (r + r' + p)^{\frac{3}{2}} - (r + r' - p)^{\frac{3}{2}}
\]
Es sei
\[
\frac{\sqrt{(r + r' + p)} + \sqrt{(r + r' - p)}}{2\sqrt{(r + r')}} = \cos 2\psi, \quad \frac{\sqrt{(r + r' + p)} - \sqrt{(r + r' - p)}}{2\sqrt{(r + r')}} = \sin\varphi\cos 2\psi
\]
[I] $\sqrt{8kT} : (r + r')^{\frac{3}{2}} = \sin\varphi : \sqrt{\cos 2\psi}$

[Ferner ist, wenn vor\"ubergehend $t = \sqrt{\frac{r + r' + p}{r + r'}}$, $u = \sqrt{\frac{r + r' - p}{r + r'}}$ gesetzt wird,]
\[
6kT = (r + r')^{\frac{3}{2}}(t - u)(tt + tu + uu) = (r + r')^{\frac{3}{2}}(t - u)(3 - \ldots)
\]
[also, wegen] $\frac{t - u}{2} = \sin\varphi$:
\[
\text{[II]} \quad \frac{\sqrt{8kT}}{(r + r')^{\frac{3}{2}}} = \sin\varphi \cdot \ldots
\]
[Nachdem $\psi$ aus Gleichung II gefunden ist, ergibt sich $p$ aus Gleichung I.]

\subsection*{[3.]}
\textbf{Zur parabolischen Bewegung.}

[Bestimmung der Elemente durch zwei Radios Vectores $r$, $r'$ und Chorde $p$.
Es sei
\begin{equation}\tag{1}
\frac{r' + r + p}{2\sqrt{rr'}} = \sin\psi, \quad \frac{r' + r - p}{2\sqrt{rr'}} = \sin\varphi.
\end{equation}

\newpage

Mit den Bezeichnungen der \textit{Theoria motus} Art.~98 und 108 hat man:
\[
\xi = \frac{r' + r + p}{2\sqrt{rr'}} = 1 + \cos f\cos F, \quad \eta = \frac{r' - r}{2\sqrt{rr'}} = \sin f\sin F,
\]
also
\[
\xi\eta = (\cos f + \cos F)\sin f\sin F
\]
oder, da
\[
\frac{p}{rr'} = \sin f\sin F,
\]
wo
\[
\xi = \frac{1 + \cos f^2 + 2\cos f\cos F}{2(\cos f + \cos F)}, \quad \eta = \frac{\sin f^2\sin F}{2(\cos f + \cos F)},
\]
so wird
\[
\frac{p}{r' + r} = \frac{\sin f\sqrt{(1 + \cos f^2 + 2\cos f\cos F)}}{1 + \cos f\cos F}, \quad \frac{r' - r}{p} = \frac{\sin F}{\sqrt{(1 + \cos f^2 + 2\cos f\cos F)}}.
\]

Hieraus:
\[
\cos f + \cos F = \frac{\cos f(1 + \cos f^2 + 2\cos f\cos F)}{1 + \cos f\cos F},
\]
\[
\cos\psi\tan\varphi = \tan f,
\]
woraus $f$ gefunden wird.

Ferner
\[
\sin\psi\sin f = \frac{\sin f(\cos f + \cos F)}{1 + \cos f\cos F},
\]
\[
\sin\psi\sin\vartheta = \frac{\sin f\sin F}{1 + \cos f\cos F},
\]
hieraus
\begin{align*}
\sin\psi\sin(F - \psi) &= \cos f\sin\psi\sin\varphi = \sin f\tan\psi\cos\psi \\
\sin\psi\cos(F - \psi) &= \sin f\sin\psi.
\end{align*}

Diese Gleichungen selbst, oder die daraus folgenden:
\[
\tan\psi\cos\psi = \tan(F - \psi)
\]
oder
\begin{align*}
\sin f\sin\psi\tan\tfrac{1}{2}\varphi &= \sin(2\psi - F) \\
\sin f\sin\psi\cot\tfrac{1}{2}\varphi &= \sin F
\end{align*}
[dienen zur Bestimmung von $F$;] die wahren Anomalien [sind dann] $F - f$
und $F + f$.

\newpage

\subsection*{[4.]}
\textbf{Verh\"altniss von Dreieck und Sector zwischen zwei Radiis vectoribus
in der parabolischen Bewegung.}

[Nach dem Vorigen findet man:
\[
m^2 - n^2 = 2(r' + r)(\cos(45^\circ - \tfrac{1}{2}\varphi)^2 - \sin(45^\circ - \tfrac{1}{2}\varphi)^2) = 2(r' + r)\sin\tfrac{1}{2}\varphi \cdot (2 + \cos\varphi),
\]
also]
\[
\frac{m^2 - n^2}{rr'} = \frac{2(r' + r)\sin\tfrac{1}{2}\varphi \cdot (2 + \cos\varphi)}{rr'}.
\]

[Hieraus ergibt sich der in der Zeit $t' - t$ beschriebene Ausschnitt]
\[
\text{Ausschnitt} = \frac{2 + \cos\varphi}{3\cos\varphi} \cdot \text{Dreieck}.
\]

Das Dreieck zwischen $r$ und $r'$ ist aber
\[
= \tfrac{1}{4}\sqrt{(r' + r + p)(r' + r - p)(r' - r + p)(-r' + r + p)} = \tfrac{1}{2}rr'\sin\psi\cos\psi,
\]
also]
\[
\frac{\text{Ausschnitt}}{\text{Dreieck}} = \frac{2 + \cos\varphi}{3\cos\varphi}.
\]

[Durch Entwickelung der Gleichung $6k(t' - t) = (r' + r + p)^{\frac{3}{2}} - (r' + r - p)^{\frac{3}{2}}$
ergibt sich:]
\[
2k(t' - t) = p\sqrt{(r' + r)} \cdot \ldots
\]

\begin{center}
\begin{tabular}{c}
Ausschnitt \\
Dreieck
\end{tabular}
$= 1 + \tfrac{1}{2}a + \tfrac{1}{2}aa + \tfrac{1}{4}a^3 + \tfrac{1}{8}a^4 + \text{etc.},$

wo $a = \ldots$ [$= \sin\psi^2$].
\end{center}

Setzt man $\frac{\text{Ausschnitt}}{\text{Dreieck}} = \beta$, so wird
\[
\beta = a - \tfrac{1}{12}aa - \tfrac{1}{24}a^3 - \tfrac{1}{45}a^4 - \tfrac{1}{40}a^5 - \text{etc.}
\]
Folglich

[Die obige Reihe im Ausdruck f\"ur $2k(t' - t)$ l\"asst sich sehr nahe dar-
stellen durch]

[Setzt man]
\[
\frac{p\sqrt{(r' + r)}}{2k(t' - t)} = x = \sin\varphi,
\]
[ist]
\[
\frac{p\sqrt{(r' + r)}}{2k(t' - t)} = \frac{1}{1 + x} - \frac{3\cos\tfrac{1}{2}\varphi}{2 + \cos\varphi}
\]
\[
\frac{(1 + x)^2 - (1 - x)^2}{(1 + x)^2} = \ldots
\]
[Zur L\"osung dieser Gleichung kann man sich eine Tafel von folgender
Form entwerfen:]

\newpage

\begin{center}
\begin{longtable}{|c|c|c|c|c|c|c|c|}
\hline
$\log x$ & $\log y$ & $\log x$ & $\log y$ & $\log x$ & $\log y$ & $\log x$ & $\log y$ \\
\hline
\endfirsthead
\hline
$\log x$ & $\log y$ & $\log x$ & $\log y$ & $\log x$ & $\log y$ & $\log x$ & $\log y$ \\
\hline
\endhead
\multicolumn{8}{r}{\textit{Fortsetzung n\"achste Seite}} \\
\endfoot
\hline
\endlastfoot
$8{,}22$ & $0{,}00000$ & $9{,}00$ & $0{,}00018$ & $9{,}32$ & $0{,}00080$ & $9{,}64$ & $0{,}00360$ \\
$8{,}23$ & $1$ & $9{,}01$ & $19$ & $9{,}33$ & $84$ & $9{,}65$ & $377$ \\
$8{,}45$ & $1$ & $9{,}02$ & $20$ & $9{,}34$ & $87$ & $9{,}66$ & $396$ \\
$8{,}46$ & $2$ & $9{,}03$ & $21$ & $9{,}35$ & $92$ & $9{,}67$ & $416$ \\
$8{,}57$ & $2$ & $9{,}04$ & $22$ & $9{,}36$ & $96$ & $9{,}68$ & $436$ \\
$8{,}58$ & $3$ & $9{,}05$ & $23$ & $9{,}37$ & $101$ & $9{,}69$ & $458$ \\
$8{,}64$ & $3$ & $9{,}06$ & $24$ & $9{,}38$ & $105$ & $9{,}70$ & $481$ \\
$8{,}65$ & $4$ & $9{,}07$ & $25$ & $9{,}39$ & $110$ & $9{,}71$ & $505$ \\
$8{,}69$ & $4$ & $9{,}08$ & $26$ & $9{,}40$ & $116$ & $9{,}72$ & $530$ \\
$8{,}70$ & $5$ & $9{,}09$ & $27$ & $9{,}41$ & $121$ & $9{,}73$ & $557$ \\
$8{,}74$ & $5$ & $9{,}10$ & $29$ & $9{,}42$ & $127$ & $9{,}74$ & $585$ \\
$8{,}75$ & $6$ & $9{,}11$ & $30$ & $9{,}43$ & $133$ & $9{,}75$ & $615$ \\
$8{,}77$ & $6$ & $9{,}12$ & $32$ & $9{,}44$ & $140$ & $9{,}76$ & $647$ \\
$8{,}78$ & $7$ & $9{,}13$ & $33$ & $9{,}45$ & $146$ & $9{,}77$ & $680$ \\
$8{,}80$ & $7$ & $9{,}14$ & $35$ & $9{,}46$ & $153$ & $9{,}78$ & $715$ \\
$8{,}81$ & $8$ & $9{,}15$ & $36$ & $9{,}47$ & $161$ & $9{,}79$ & $752$ \\
$8{,}83$ & $8$ & $9{,}16$ & $38$ & $9{,}48$ & $168$ & $9{,}80$ & $791$ \\
$8{,}84$ & $9$ & $9{,}17$ & $40$ & $9{,}49$ & $176$ & $9{,}81$ & $833$ \\
$8{,}85$ & $9$ & $9{,}18$ & $42$ & $9{,}50$ & $185$ & $9{,}82$ & $877$ \\
$8{,}86$ & $10$ & $9{,}19$ & $44$ & $9{,}51$ & $194$ & $9{,}83$ & $924$ \\
$8{,}88$ & $10$ & $9{,}20$ & $46$ & $9{,}52$ & $203$ & $9{,}84$ & $973$ \\
$8{,}89$ & $11$ & $9{,}21$ & $48$ & $9{,}53$ & $213$ & $9{,}85$ & $1026$ \\
$8{,}90$ & $11$ & $9{,}22$ & $50$ & $9{,}54$ & $223$ & $9{,}86$ & $1082$ \\
$8{,}91$ & $12$ & $9{,}23$ & $53$ & $9{,}55$ & $234$ & $9{,}87$ & $1142$ \\
$8{,}92$ & $13$ & $9{,}24$ & $55$ & $9{,}56$ & $245$ & $9{,}88$ & $1205$ \\
$8{,}93$ & $13$ & $9{,}25$ & $58$ & $9{,}57$ & $257$ & $9{,}89$ & $1273$ \\
$8{,}94$ & $14$ & $9{,}26$ & $60$ & $9{,}58$ & $270$ & $9{,}90$ & $1346$ \\
$8{,}95$ & $14$ & $9{,}27$ & $63$ & $9{,}59$ & $283$ & & \\
$8{,}96$ & $15$ & $9{,}28$ & $66$ & $9{,}60$ & $297$ & & \\
$8{,}97$ & $16$ & $9{,}29$ & $69$ & $9{,}61$ & $311$ & & \\
$8{,}98$ & $17$ & $9{,}30$ & $73$ & $9{,}62$ & $327$ & & \\
$8{,}99$ & $0{,}00017$ & $9{,}31$ & $0{,}00076$ & $9{,}63$ & $0{,}00343$ & $0{,}00$ & $0{,}02558$ \\
\hline
\end{longtable}
\end{center}

\begin{flushright}
$42^*$
\end{flushright}

\newpage

\subsection*{[5.]}
\textbf{Gauss an Olbers, G\"ottingen. 7.~Januar 1815.}

Ich habe Ihnen in meinem vorigen Briefe die Ann\"aherungsformel f\"ur die
Zwischenzeit
\[
t'' - t = 3k\sqrt{(r + r'')} \cdot \sqrt{(1 - \ldots)} \cdot u^3
\]
geschrieben und gesagt, dass sie in den meisten F\"allen hinreichend genau sei.
Ich h\"atte sagen sollen, es werde in der Praxis kein Fall vorkommen, wo sie
nicht \"uberfl\"ussig genau w\"are. So lange $k < \tfrac{1}{2}(r + r'')$, afficirt der Fehler noch
nicht die $5^\text{te}$ Decimale. Es ist eine Lust, zu sehen, wie bequem die Versuche
f\"ur $u$ danach durchgef\"uhrt werden. Jeder Versuch erfordert bloss 8 Auf-
schlagungen, 3 in den Logarithmen, 4 in den Sinustafeln und eine in meiner
H\"ulfstafel f\"ur Logarithmen von Summen und Differenzen. Um so wenig wie
m\"oglich zu schreiben zu haben, setze ich $r^{\frac{1}{2}} = s$, $r''^{\frac{1}{2}} = s''$; da-
durch wird, wenn man noch $\frac{k}{s + s''} = \sin Q$ setzt:
\[
t'' - t = k\sqrt{(s + s'')} \cdot \ldots
\]
Die ganze Rechnung besteht also vorher nur noch in folgendem:
\begin{center}
\begin{tabular}{c|c}
$\sqrt{r}$ & $\sqrt{r''}$ \\
\hline
$k$ & 
\end{tabular}
\end{center}

\begin{flushright}
[$^*$ Hier bezeichnet, abweichend von den vorstehenden Notizen, $k$ die Chorde, $\frac{k}{\ldots}$ die im Vorigen
mit $k$ bezeichnete Constante, vgl.~S.~330.]
\end{flushright}

\newpage

\subsection*{[6.]}
\textbf{Gauss an Olbers, G\"ottingen, 13.~Januar 1815.}

Ich habe mich diese Woche hindurch noch mit der Theorie der Cometen-
bahnen besch\"aftigt, und bin noch auf eine Vereinfachung gekommen, nach
welcher mir nun nichts mehr zu w\"unschen \"ubrig zu bleiben scheint. Ich
finde nemlich f\"ur die Zwischenzeit den Ausdruck:
\[
t'' - t = \frac{k\sqrt{(r + r'')}}{\cos Q'} \cdot \sqrt{(1 + \ldots)}
\]
der in allen in der Praxis vorkommenden F\"allen genau genug ist. Nach
v\"olliger Strenge sollte die letzte Quadratwurzel eine unendliche Reihe sein,
die, wenn ich der K\"urze halber $\frac{k^2}{rr''} = x$ setze, nach meiner Entwickelung
wird $= \sqrt{(1 - x - 2x^2 - 8x^3 - \ldots)}$ und wo also gl\"ucklicherweise das zweite
Glied fehlt. Selbst, wenn $k = \sqrt{(r + r'')}$, w\"urde $x$ nur etwa $\tfrac{1}{4}$, also der Fehler,
wenn man sich auf $\sqrt{(1 - x)}$ einschr\"ankt, nur etwa $\tfrac{1}{100}$ des Ganzen [$^*$].

Ich wende nun obige Formel auf folgende Art an. Ich setze
\[
u = A\tan Q,
\]
wodurch
\[
k = \frac{A}{\cos Q} \quad \text{und} \quad p = A\cos\psi \ldots
\]
Man hat also
\[
\frac{t'' - t}{A\sqrt{(r + r'')}} = \ldots
\]
oder, indem ich folgende Bezeichnungen einf\"uhre
\begin{align*}
\frac{mmAA}{cc(t'' - t)^2} &= \delta, & \frac{mmAA}{cc(t'' - t)^2} &= \delta'', \\
\frac{\delta B}{b} &= \zeta, & \frac{\delta''B''}{b''} &= \zeta'', \\
\frac{cc}{4(t'' - t)^2} &= \mu, & \frac{c''c''}{4(t'' - t)^2} &= \mu'',
\end{align*}
wird:
\[
\frac{t'' - t}{A\sqrt{(r + r'')}} = \delta\sqrt{(1 + \zeta\tan Q + \mu\zeta^2)} + \delta''\sqrt{(1 + \zeta''\tan Q + \mu''\zeta''^2)}.
\]

\newpage

Etwas einfacheres ist wohl nicht zu w\"unschen. Man k\"onnte noch zwei
H\"ulfswinkel $\vartheta$, $\vartheta''$ einf\"uhren, indem man
\begin{align*}
\delta\zeta &= \lambda\sin\vartheta, & \delta''\zeta'' &= \lambda''\sin\vartheta'' \\
\delta &= \lambda\cos\vartheta, & \delta'' &= \lambda''\cos\vartheta''
\end{align*}
setzte, wodurch die Gleichung w\"urde
\[
\frac{t'' - t}{A\sqrt{(r + r'')}} = \sqrt{(\delta\delta\cos Q^2 + \lambda\lambda\cos Q - \ldots)} + \sqrt{(\delta''\delta''\cos Q^2 + \lambda''\lambda''\cos Q - \ldots)}.
\]
Man k\"onnte selbst noch weiter gehen und durch bekannte Verwandlungen
der Gleichung folgende Form geben:
\[
\frac{t'' - t}{A\sqrt{(r + r'')}} = \tfrac{1}{2}\mu\sqrt{(1 + \nu'\nu'\cos Q + \zeta'')} + \tfrac{1}{2}\mu''\sqrt{(1 + \nu''\nu''\cos Q + \zeta''')}.
\]

\begin{flushright}
[$^*$ In einem Exemplar der \textit{Abhandlung "Observationes Cometae secundi A.~1813 etc."} findet
sich die handschriftliche Bemerkung:] NB. Die Reihe $1 - x - 2x^2 - 8x^3 - 45x^4 - 172x^5 \ldots$
convergirt, so lange $x < \tfrac{1}{4}$; f\"ur diese Grenze, der $k = r + r''$ entspricht, wird ihr Werth $= \tfrac{1}{2}$.
\end{flushright}

\newpage

Allein der Gewinn[$^*$] w\"urde zu unerheblich sein, um die M\"uhe dieser
Verwandlungen zu belohnen. Immer ist die Gleichung vollkommen strenge,
sobald man statt des Nenners im ersten Theile die Reihe setzt
\[
1 + \frac{2\mu'}{\cos Q'^2} + \frac{8\mu'^2}{\cos Q'^4} + \ldots
\]
Das Problem ist also auf Eine Gleichung mit sieben bekannten Gr\"ossen
gebracht, und ich glaube nicht, dass es m\"oglich ist, es auf eine geringere
Anzahl zu reduciren.

Die allgemeine Aufgabe, nemlich, wenn der mittelste Ort unvollst\"andig
bekannt ist, und man bloss irgend einen gr\"ossten Kreis kennt, worin er liegt
--- von welcher das Gegenw\"artige eigentlich nur der einzelne Fall ist, wo
dieser gr\"osste Kreis durch den mittlern Ort geht --- zu welcher man seine
Zuflucht nehmen muss, so oft dieser Ort nahe bei der Richtung der geo-
centrischen Bewegung liegt, diese allgemeine Aufgabe habe ich mir nun auch
auf eine analoge Art durchgef\"uhrt und ein Musterbeispiel sorgf\"altig be-
rechnet [$^*$]. Der Natur der Sache nach ist die Arbeit hier viel verwickelter,
ich komme am Ende auf zwei Gleichungen mit $u$[$^{**}$] bekannten Gr\"ossen, und
zweifle auch, ob es m\"oglich ist, diese Anzahl kleiner zu machen. Wenn nach
allen Abk\"urzungen jetzt die Berechnung einer Cometenbahn im ersten Fall
im Minimum 1 Stunde erfordert (so ist es etwa bei mir), so m\"ochte der an-
dere Fall leicht 2--3 Stunden erfordern.

\begin{flushright}
[$^*$) Er w\"urde nemlich bloss darin bestehen, dass man sogleich das Maximum und Minimum des
Werthes der beiden Wurzelgr\"ossen beurtheilen k\"onnte, wodurch man, anfangs $u$ vernachl\"assigend, Grenzen
f\"ur $\cos Q'$ und folglich f\"ur $Q$ erhielte. Bei der Berechnung jedes einzelnen Versuchs werden doch immer
7 Aufschlagungen erfordert, man w\"ahle welche Form man wolle.]

[$^*$) Siehe die folgende Notiz.]

[$^{**}$) Jeder Versuch 11 Aufschlagungen erfordernd.]
\end{flushright}

\newpage

\subsection*{[7.]}
\textbf{Allgemeine Theorie der Berechnung der Cometenbahnen}
[nebst Anwendung auf den zweiten Cometen des Jahrs 1813].

\begin{center}
\begin{tabular}{|c|c|c|c|c|}
\hline
\multicolumn{2}{|c|}{Beobachtungszeiten} & Beobachtete L\"angen & Beobachtete Breiten \\
\hline
$t$ & $[1813$ April] $7{,}55002$ & $\alpha = 271^\circ 16' 38''$ & $\beta = +29^\circ 2' 0''$ \\
$t'$ & $\quad\,,\quad\,,\quad 14{,}54694$ & $\alpha' = 266^\circ 27' 22''$ & $\beta' = -22^\circ 52' 18''$ \\
$t''$ & $\quad\,,\quad\,,\quad 21{,}59931$ & $\alpha'' = 256^\circ 48' 8''$ & $\beta'' = +9^\circ 53' 12''$ \\
\hline
\multicolumn{2}{|c|}{L\"angen der Sonne} & \multicolumn{2}{c|}{Logarithmen der Abst\"ande} \\
\hline
$\Theta = 17^\circ 47' 41''$ & $\log R = 0{,}00091$ \\
$\Theta' = 24^\circ 38' 45''$ & $\log R' = 0{,}00175$ \\
$\Theta'' = 31^\circ 31' 25''$ & $\log R'' = 0{,}00260$ \\
\hline
\end{tabular}
\end{center}

[$r$, $r'$, $r''$ die Abst\"ande von der Sonne,
$p$, $p'$, $p''$ die Abst\"ande von der Erde.]

\textbf{Vorl\"aufige Operationen:}
\begin{align*}
pp - 2pR\cos\beta\cos(\alpha - \Theta) + RR &= r^2 \\
p''p'' - 2p''R''\cos\beta''\cos(\alpha'' - \Theta'') + R''R'' &= r''^2
\end{align*}
[Es setzt man, wie in der Abhandlung \textit{Observationes cometae secundi A.~MDCCCXIII etc.}]
\begin{align*}
\cos\psi &= \cos\beta\cos(\alpha - \Theta), & \cos\psi'' &= \cos\beta''\cos(\alpha'' - \Theta'') \\
R\sin\psi &= B, & R''\sin\psi'' &= B''
\end{align*}
[so wird]

\newpage

Legt man durch den ersten und dritten geocentrischen Ort einen gr\"ossten
Kreis, der die Ekliptik im Punkte $N$ unter dem Neigungswinkel $J$ schneide,
so ist, wenn die L\"ange dieses Punkts ebenfalls mit $N$ bezeichnet wird,
\begin{align*}
rr &= (p - R\cos\psi)^2 + BB \\
r''r'' &= (p'' - R''\cos\psi'')^2 + B''B''.
\end{align*}

\begin{align*}
\tan J\sin(\alpha - N) &= \tan\beta \\
\tan J\sin(\alpha'' - N) &= \tan\beta'';
\end{align*}
$J$ und $N$ k\"onnen aus den Formeln, vgl.~\textit{Theoria motus corporum coelestium}
Art.~78,
\begin{align*}
\tan J\sin(\alpha - N) &= \tan\beta \\
\tan J\cos(\alpha - N) &= \frac{\tan\beta'' - \tan\beta\cos(\alpha - \alpha'')}{\sin(\alpha'' - \alpha)}
\end{align*}
berechnet werden.

Indem man die L\"angen vom Punkte $N$ aus z\"ahlt, hat man f\"ur die Chorde
zwischen dem ersten und dritten Ort
\begin{multline*}
\text{Chorde}^2 = \\
\{p''\cos\beta''\cos(\alpha'' - N) - p\cos\beta\cos(\alpha - N) - R''\cos(\Theta'' - N) + R\cos(\Theta - N)\}^2 \\
+ \{p''\cos\beta''\sin(\alpha'' - N) - p\cos\beta\sin(\alpha - N) - R''\sin(\Theta'' - N) + R\sin(\Theta - N)\}^2 \\
+ \{p''\sin\beta'' - p\sin\beta\}^2.
\end{multline*}

Man setze
\begin{align*}
R''\cos(\Theta'' - N) - R\cos(\Theta - N) &= f\cos(F - N) \\
R''\sin(\Theta'' - N) - R\sin(\Theta - N) &= f\sin(F - N)
\end{align*}
und bezeichne mit $K$, $K''$ die Abst\"ande des ersten und dritten geocentrischen
Orts vom Punkte $N$; alsdann ist
\begin{align*}
\cos\beta\cos(\alpha - N) &= \cos K, & \cos\beta\sin(\alpha - N) &= \cos J\sin K, & \sin\beta &= \sin J\sin K \\
\cos\beta''\cos(\alpha'' - N) &= \cos K'', & \cos\beta''\sin(\alpha'' - N) &= \cos J\sin K'', & \sin\beta'' &= \sin J\sin K'',
\end{align*}
womit wird
\begin{multline*}
\text{Chorde}^2 = (p''\cos K'' - p\cos K - f\cos(F - N))^2 \\
+ (p''\cos J\sin K'' - p\cos J\sin K - f\sin(F - N))^2 + (p''\sin J\sin K'' - p\sin J\sin K)^2.
\end{multline*}

Man kann die Gr\"ossen $f$, $F$, $K$, $K''$ aus den Formeln
\begin{align*}
f\cos(F - \Theta) &= R''\cos(\Theta'' - \Theta) - R \\
f\sin(F - \Theta) &= R''\sin(\Theta'' - \Theta) \\
\tan K &= \frac{B}{R\cos\psi} \\
\tan K'' &= \frac{B''}{R''\cos\psi''}
\end{align*}
berechnen und noch setzen
\begin{align*}
f\cos J\sin(F - N) &= g\sin G \\
f\cos J\cos(F - N) &= g\cos G.
\end{align*}

\newpage

[Legt man nun durch den mittlern geocentrischen Ort einen gr\"ossten
Kreis, welcher den durch die beiden \"ausseren \"Orter gehenden gr\"ossten Kreis
rechtwinklig schneidet und nennt man $Q$ die L\"ange des Punkts, in dem
dieser Kreis die Ekliptik schneidet, so ist]
\begin{equation}\tag{VI}
\tan(Q - N) = \tan J \cdot \frac{\ldots}{\ldots}
\end{equation}
\begin{equation}
\tan J \cos(\alpha' - N) = \tan\beta' \cos(Q - N) + \ldots
\end{equation}

[In unserm Beispiel wird]
\begin{center}
\begin{tabular}{ll}
$\cos\beta \ldots\ldots 9{,}94168$ & \\
$\cos\psi \ldots\ldots 9{,}99350$ & \\
$\cos(\alpha'' - \Theta'') \ldots 9{,}84736$ & \\
$\cos(\alpha - \Theta) \ldots\ldots 9{,}45379$ & \\
$\cos\psi \ldots\ldots 9{,}39547$ & \\
$\cos(\alpha'' - \alpha) \ldots\ldots 9{,}84086$ & \\
$R \ldots\ldots 0{,}00091$ & \\
$R'' \ldots\ldots 0{,}00260$ & \\
$\sin\psi \ldots\ldots 9{,}98615$ & \\
$\sin\psi'' \ldots\ldots 9{,}85778$ & \\
$A \ldots\ldots 9{,}98706$ & $A'' \ldots\ldots 9{,}86038$ \\
$R'' \ldots\ldots 0{,}00260$ & \\
$\tan\beta \ldots\ldots 9{,}74435$ & $N = 200^\circ 22' 6''$ \\
$\cos(\Theta'' - \Theta) \ldots\ldots 9{,}98741$ & \\
$\cos(\alpha'' - \alpha) \ldots\ldots 9{,}98599$ & \\
$\alpha' - N = 16^\circ 5' 16''$ & \\
$\sin(\Theta'' - \Theta) \ldots\ldots 9{,}37535$ & \\
$[R''\cos(\Theta'' - \Theta)] \ldots\ldots 9{,}99001$ & \\
$\tan\beta'' \ldots\ldots 9{,}24127$ & \\
$F - N = 223^\circ 21' 51''$ & \\
$-R \ldots\ldots 0{,}00091$ & \\
$9{,}56010$ & \\
$f\cos(F - \Theta) = 8{,}39501$ & \\
$f\sin(F - \Theta) = 9{,}37795$ & \\
$\sin(\alpha'' - \alpha) \ldots\ldots 9{,}39786$ & \\
$\tan J \cos(\alpha' - N) \ldots\ldots 0{,}16224$ & \\
$\sin(F - \Theta) \ldots\ldots 9{,}99766$ & \\
$\tan J \sin(\alpha' - N) \ldots\ldots 9{,}74435$ & \\
$f = 9{,}38029$ & \\
$\cos(\alpha - N) \ldots\ldots 9{,}97041$ & \\
$F - Q = 95^\circ 56' 16''$ & \\
$\tan J \ldots\ldots 0{,}19183$ & \\
$Q = 17^\circ 47' 41''$ & \\
$\sin J \ldots\ldots 9{,}92487$ & \\
$F = 113^\circ 43' 57''$ & $\cos J = 9{,}73304$ \\
$\alpha - N = 20^\circ 54' 32''$ & \\
\end{tabular}
\end{center}

\newpage

[berechnen und noch setzen]
\begin{align*}
f\cos J\sin(F - N) &= g\sin G \\
f\cos J\cos(F - N) &= g\cos G.
\end{align*}

[Legt man nun durch den mittlern geocentrischen Ort einen gr\"ossten
Kreis, welcher den durch die beiden \"ausseren \"Orter gehenden gr\"ossten Kreis
rechtwinklig schneidet und nennt man $Q$ die L\"ange des Punkts, in dem
dieser Kreis die Ekliptik schneidet, so ist]
\begin{equation}\tag{VI}
\tan(Q - N) = \tan J \cdot \frac{\sin K''\cos K \sin(\alpha' - N) - \sin K\cos K''\sin(\alpha'' - N)}{\cos K''\cos K \sin(\alpha'' - \alpha)}
\end{equation}
\begin{equation}
\tan J \cos(\alpha' - N) = \tan\beta' \cos(Q - N) + \frac{\sin(Q - N)}{\sin K''\cos K \sin(\alpha'' - \alpha)} \cdot \ldots
\end{equation}

[In unserm Beispiel wird]
\begin{center}
\begin{tabular}{ll}
$\cos\beta \ldots\ldots 9{,}94168$ & \\
$\cos\beta'' \ldots\ldots 9{,}99350$ & \\
$\cos(\alpha'' - \Theta'') \ldots 9{,}84736$ & \\
$\cos(\alpha - \Theta) \ldots\ldots 9{,}45379$ & \\
$\cos\psi \ldots\ldots 9{,}39547$ & \\
$\cos(\alpha'' - \alpha) \ldots\ldots 9{,}84086$ & \\
$R \ldots\ldots 0{,}00091$ & \\
$R'' \ldots\ldots 0{,}00260$ & \\
$\sin\psi \ldots\ldots 9{,}98615$ & \\
$\sin\psi'' \ldots\ldots 9{,}85778$ & \\
$A \ldots\ldots 9{,}98706$ & $A'' \ldots\ldots 9{,}86038$ \\
$R'' \ldots\ldots 0{,}00260$ & \\
$\tan\beta \ldots\ldots 9{,}74435$ & $N = 200^\circ 22' 6''$ \\
$\cos(\Theta'' - \Theta) \ldots\ldots 9{,}98741$ & \\
$\cos(\alpha'' - \alpha) \ldots\ldots 9{,}98599$ & \\
$\alpha' - N = 16^\circ 5' 16''$ & \\
$\sin(\Theta'' - \Theta) \ldots\ldots 9{,}37535$ & \\
$[R''\cos(\Theta'' - \Theta)] \ldots\ldots 9{,}99001$ & \\
$\tan\beta'' \ldots\ldots 9{,}24127$ & \\
$F - N = 223^\circ 21' 51''$ & \\
$-R \ldots\ldots 0{,}00091$ & \\
$9{,}56010$ & \\
$f\cos(F - \Theta) \ldots\ldots 8{,}39501$ & \\
$f\sin(F - \Theta) \ldots\ldots 9{,}37795$ & \\
$\sin(\alpha'' - \alpha) \ldots\ldots 9{,}39786$ & \\
$\tan J \cos(\alpha' - N) \ldots\ldots 0{,}16224$ & \\
$\sin(F - \Theta) \ldots\ldots 9{,}99766$ & \\
$\tan J \sin(\alpha' - N) \ldots\ldots 9{,}74435$ & \\
$f \ldots\ldots 9{,}38029$ & \\
$\cos(\alpha - N) \ldots\ldots 9{,}97041$ & \\
$F - Q = 95^\circ 56' 16''$ & \\
$\tan J \ldots\ldots 0{,}19183$ & \\
$Q = 17^\circ 47' 41''$ & \\
$\sin J \ldots\ldots 9{,}92487$ & \\
$F = 113^\circ 43' 57''$ & $\cos J \ldots\ldots 9{,}73304$ \\
$\alpha - N = 20^\circ 54' 32''$ & \\
\end{tabular}
\end{center}

\newpage

\section*{Zur elliptischen Bewegung}
\addcontentsline{toc}{section}{Zur elliptischen Bewegung}

\subsection*{[1.]}

$M$ mittlere Anomalie, $E$ excentrische Anomalie, $e$ Excentricit\"at, $a$ halbe
 grosse Axe, $r$ Radius Vector, $\pi$ Perihelium, $\lambda$ L\"ange in der Bahn, $k = \sqrt{M}$.

\begin{align*}
r\cos\lambda &= a(\cos E - e) & r\sin\lambda &= a\sin E \sqrt{(1 - ee)} \\
r &= a(1 - e\cos E) & \tan\tfrac{1}{2}\lambda &= \sqrt{\tfrac{1 + e}{1 - e}} \tan\tfrac{1}{2}E \\
M &= E - e\sin E & \frac{dr}{dE} &= ae\sin E \\
\frac{dM}{dE} &= r & \frac{d\lambda}{dE} &= \frac{a\sqrt{(1 - ee)}}{r}
\end{align*}

\newpage

\section*{Zur parabolischen Bewegung}
\addcontentsline{toc}{section}{Zur parabolischen Bewegung}

\subsection*{[1.]}

$k = \sqrt{M}$, $q$ Periheldistanz, $r$ Radius Vector, $v$ wahre Anomalie, $t$ Zeit
seit dem Periheldurchgang.

\begin{align*}
r &= \frac{q}{\cos\tfrac{1}{2}v^2} = \frac{2q}{1 + \cos v} \\
\tan\tfrac{1}{2}v^3 + 3\tan\tfrac{1}{2}v &= \frac{6k(t - \tau)}{q\sqrt{2q}} = W \\
r\cos v &= q(1 - \tan\tfrac{1}{2}v^2) \\
r\sin v &= 2q\tan\tfrac{1}{2}v
\end{align*}

\begin{align*}
mm &= 2q \quad nn = 2q\tan\tfrac{1}{2}v^2 \quad m + n = 2\sqrt{q}\cos\tfrac{1}{2}v \\
m - n &= 2\sqrt{q}\sin\tfrac{1}{2}v \quad mn = 2q\tan\tfrac{1}{2}v
\end{align*}

\begin{align*}
W &= \frac{6k(t - \tau)}{(2q)^{\frac{3}{2}}} = \frac{3k(t - \tau)}{mm\sqrt{m}} \\
\frac{W}{3} &= \frac{k(t - \tau)}{m\sqrt{mm}} = \tan\tfrac{1}{2}v + \tfrac{1}{3}\tan\tfrac{1}{2}v^3
\end{align*}

\begin{align*}
t' - t &= \frac{mm'm' - nnn'}{6kk'} \cdot \ldots \\
&= \frac{1}{6kk'}(r' + r + p)^{\frac{3}{2}} - (r' + r - p)^{\frac{3}{2}}
\end{align*}

\newpage

\end{document}
